\documentclass[12pt,oneside]{report}

\usepackage[a4paper,left=1.5in,right=1in,top=1in,bottom=1in]{geometry}
\usepackage[T1]{fontenc}
\usepackage[utf8]{inputenc}
\usepackage{newtxtext}
\usepackage{setspace}
\usepackage{ragged2e}
\usepackage{titlesec}
\usepackage{graphicx}
\usepackage{booktabs}
\usepackage{array}
\usepackage{longtable}
\usepackage{multirow}
\usepackage{amsmath,amssymb,amsfonts}
\usepackage{newtxmath}
\usepackage{algorithm}
\usepackage{algorithmic}
\usepackage{listings}
\usepackage{xcolor}
\usepackage{hyperref}
\usepackage{enumitem}
\usepackage{caption}
\usepackage{float}
\usepackage{fancyhdr}

\hypersetup{colorlinks=true,linkcolor=black,citecolor=black,urlcolor=black}
\newcommand{\safeincludegraphics}[2][]{%
  \IfFileExists{#2}{%
    \includegraphics[#1]{#2}%
  }{%
    \fbox{\parbox[c][2.0in][c]{0.82\textwidth}{\centering Missing figure file:\\\texttt{\detokenize{#2}}}}%
  }%
}

\titleformat{\chapter}[display]
  {\normalfont\bfseries\fontsize{16}{19}\selectfont\centering}
  {\chaptertitlename\ \thechapter}{10pt}{\fontsize{16}{19}\selectfont}
\titleformat{\section}
  {\normalfont\bfseries\fontsize{14}{17}\selectfont}
  {\thesection}{1em}{}
\titleformat{\subsection}
  {\normalfont\bfseries\fontsize{12}{15}\selectfont}
  {\thesubsection}{1em}{}
\captionsetup{font=small,labelfont=bf}
\setlength{\parindent}{0.5in}
\setlength{\parskip}{0pt}
\justifying
\onehalfspacing

\lstdefinestyle{codestyle}{
    basicstyle=\ttfamily\footnotesize,
    breaklines=true,
    columns=fullflexible,
    keepspaces=true,
    frame=single,
    numbers=left,
    numberstyle=\tiny,
    keywordstyle=\bfseries,
    commentstyle=\itshape,
    showstringspaces=false,
    tabsize=4
}
\lstset{style=codestyle}

\pagestyle{fancy}
\fancyhf{}
\fancyfoot[C]{\thepage}
\renewcommand{\headrulewidth}{0pt}

\begin{document}

\pagenumbering{roman}

\begin{titlepage}
\centering
\vspace*{0.3in}
{\fontsize{16}{20}\selectfont\bfseries Ontology-Driven Self-Healing Defense Mechanisms for Graph Neural Networks Against Advanced Poisoning Attacks\par}
\vspace{0.35in}
{\fontsize{14}{17}\selectfont\bfseries M.Tech Thesis Project Report\par}
\vspace{0.25in}
{\fontsize{12}{15}\selectfont Submitted in partial fulfillment of the requirements for the award of the degree of\par}
\vspace{0.18in}
{\fontsize{14}{17}\selectfont\bfseries Master of Technology\par}
\vspace{0.18in}
{\fontsize{12}{15}\selectfont in Computer Science and Engineering / Information Technology / Artificial Intelligence\par}
\vspace{0.45in}
{\fontsize{12}{15}\selectfont\bfseries Submitted by\par}
\vspace{0.1in}
{\fontsize{14}{17}\selectfont\bfseries [Student Name]\par}
{\fontsize{12}{15}\selectfont Roll Number: [Roll Number]\par}
\vspace{0.35in}
{\fontsize{12}{15}\selectfont\bfseries Under the Guidance of\par}
\vspace{0.1in}
{\fontsize{14}{17}\selectfont\bfseries [Advisor Name]\par}
{\fontsize{12}{15}\selectfont [Designation]\par}
\vspace{0.35in}
\fbox{\parbox[c][1.0in][c]{1.4in}{\centering University\\Logo}}
\vfill
{\fontsize{12}{15}\selectfont Department of [Department Name]\par}
{\fontsize{12}{15}\selectfont [Institution Name]\par}
{\fontsize{12}{15}\selectfont [University Name]\par}
{\fontsize{12}{15}\selectfont May 2026\par}
\end{titlepage}

\chapter*{Certificate}
\addcontentsline{toc}{chapter}{Certificate}
This is to certify that the thesis project report entitled \textbf{``Ontology-Driven Self-Healing Defense Mechanisms for Graph Neural Networks Against Advanced Poisoning Attacks''} submitted by \textbf{[Student Name]}, Roll Number \textbf{[Roll Number]}, is a bonafide record of the project work carried out under my supervision in partial fulfillment of the requirements for the award of the degree of Master of Technology. The work reported in this thesis has been developed using a local JAX/Flax Graph Neural Network framework containing modules for graph datasets, GCN/GAT models, adversarial attack generation, STRUC-GUARD+ defense, ontology-driven self-healing, advanced robustness metrics, and visualization.

The thesis work investigates the vulnerability of message-passing graph neural networks under structural poisoning and feature evasion attacks, and proposes a defense framework that combines semantic reasoning, topology-aware filtering, small-world reconstruction, and adversarial retraining. To the best of my knowledge, the work has not been submitted elsewhere in the same form for the award of any degree or diploma.

\vspace{0.8in}
\noindent Guide Signature:\hfill Head of Department:\hspace{0.6in}\\[0.4in]
\noindent Date:\hfill Place:\hspace{1.3in}

\chapter*{Declaration}
\addcontentsline{toc}{chapter}{Declaration}
I hereby declare that this thesis project report entitled \textbf{``Ontology-Driven Self-Healing Defense Mechanisms for Graph Neural Networks Against Advanced Poisoning Attacks''} is the result of my own work carried out using the local JAX/Flax GNN implementation. All research papers, datasets, software libraries, algorithms, and implementation references used in this work have been acknowledged through appropriate citation. The implementation and result analysis presented in this document are based on the local project directory, including the \texttt{attacks/}, \texttt{defense/}, \texttt{evaluation/}, \texttt{models/}, \texttt{datasets/}, and \texttt{visualization/} modules.

\vspace{0.8in}
\noindent Student Signature:\hfill Date:\hspace{1.2in}

\chapter*{Acknowledgements}
\addcontentsline{toc}{chapter}{Acknowledgements}
I express my sincere gratitude to my project guide, faculty members, laboratory staff, and the Department of [Department Name] for their continuous academic and technical support. This project required the combined study of graph representation learning, adversarial machine learning, network security, financial fraud analytics, temporal graph processing, and Python-based scientific computing. The guidance received throughout the project helped transform an initial attack-defense prototype into a rigorous thesis framework built around verifiable adversarial degradation and recovery.

I acknowledge the open-source communities behind JAX, Flax, Optax, NumPy, SciPy, NetworkX, scikit-learn, Matplotlib, and PyTorch Geometric. I also acknowledge the authors of GNNGuard, Nettack, Metattack, topology attack research, and graph adversarial learning surveys, whose work provided the theoretical foundation for this project. Finally, I thank my family, classmates, and friends for their encouragement and support during the implementation, evaluation, and documentation of this thesis.

\chapter*{Abstract}
\addcontentsline{toc}{chapter}{Abstract}
Graph Neural Networks (GNNs) have become central to learning over non-Euclidean data such as citation networks, transaction graphs, social graphs, and cybersecurity dependency graphs. Their predictive strength arises from message passing: each node representation is updated by aggregating feature information from its graph neighborhood. The same mechanism creates an adversarial surface. Localized changes in edges or features can propagate through normalized adjacency and corrupt node representations far beyond the directly perturbed node. This thesis develops a JAX/Flax adversarial robustness framework for Graph Convolutional Networks (GCNs), using Cora and Elliptic Bitcoin datasets as static and temporal graph benchmarks.

The project originated from a practical weakness in the initial experimental setting: baseline poisoning attacks such as Nettack and Meta Attack produced only marginal accuracy degradation, approximately two percentage points in early result tables. Such attacks are insufficient for evaluating defense resilience. The framework therefore introduces attack strengthening through validation-gated calibration, high-confidence target selection, surrogate gradient-matching, and structural bottleneck targeting using edge betweenness centrality. DICE and Edge Flip are redesigned to damage homophilous graph structure, while Elliptic-specific temporal perturbation breaks consistency across 49 transaction snapshots.

The proposed defense, STRUC-GUARD+ with ontology-driven self-healing, extends the base-paper GNNGuard architecture by moving beyond local cosine similarity. It constructs semantic topic sets, flags suspicious edges through topic Jaccard mismatch, prunes high-centrality low-similarity bridges, applies Bose-Einstein-inspired degree consistency checks, performs residual feature denoising, reconstructs only small-world-safe KNN edges, detects temporal drift in Elliptic snapshots, and retrains adversarially on the defended graph. The evaluation extends accuracy and F1-score with Attack Success Rate, Neighborhood Entropy, Embedding Drift, Homophily Drop, Bose-Einstein Fitness, Assortativity Coefficient, Clean Label Recovery, and Injected Edge Pruning. Final gated results show all listed Cora and Elliptic attacks exceeding the configured 50\% baseline degradation gate, with defended accuracy returning to or exceeding the clean baseline.

\clearpage
\tableofcontents
\clearpage
\listoftables
\clearpage
\listoffigures
\clearpage

\chapter*{List of Abbreviations}
\addcontentsline{toc}{chapter}{List of Abbreviations}
\begin{longtable}{p{1.3in}p{4.3in}}
APL & Average Path Length\\
ASR & Attack Success Rate\\
BE & Bose-Einstein\\
CLR & Clean Label Recovery\\
DICE & Delete Internally, Connect Externally\\
GCN & Graph Convolutional Network\\
GAT & Graph Attention Network\\
GNN & Graph Neural Network\\
JAX & Just After eXecution, a composable automatic differentiation framework\\
KNN & k-Nearest Neighbors\\
PGD & Projected Gradient Descent\\
SFIE & Scale-Free Integrity Engine\\
STRUC-GUARD+ & Enhanced Structural Defense with Ontology-Driven Self-Healing\\
\end{longtable}

\clearpage
\pagenumbering{arabic}

% ==========================================================
\chapter{Introduction}
% ==========================================================

\section{Learning on Non-Euclidean Data}
The development of graph neural networks marks a significant shift from Euclidean learning over vectors, images, and sequences to learning over relational structures. In a conventional image model, pixels are arranged on a regular grid. In a text model, tokens follow a sequence. In a graph, however, each entity is embedded in an irregular topology. The data object is not merely a matrix of attributes but a relational system:

$$
G=(V,E,X),
$$

where $V$ is the set of nodes, $E$ is the set of edges, and $X\in\mathbb{R}^{|V|\times d}$ is the node feature matrix. The graph may represent papers and citations, accounts and financial transactions, hosts and network flows, users and social relations, or molecules and bonds. The essential computational challenge is that each node's label often depends on its own attributes and on the attributes of neighboring nodes. A GNN formalizes this dependence through message passing.

The Graph Convolutional Network (GCN) used in this project follows the normalized propagation rule:

$$
H^{(\ell+1)}=\sigma\left(\hat{A}H^{(\ell)}W^{(\ell)}\right),
$$

where $H^{(0)}=X$, $W^{(\ell)}$ is a trainable matrix, $\sigma$ is a nonlinear activation, and

$$
\hat{A}=\tilde{D}^{-\frac{1}{2}}\tilde{A}\tilde{D}^{-\frac{1}{2}},\qquad \tilde{A}=A+I.
$$

This equation is compact but security-critical. The same adjacency operator that enables relational learning also acts as a channel through which adversarial perturbations can spread. If an attacker changes a small set of edges or feature values, the perturbation may enter the hidden state through $\hat{A}XW^{(0)}$, propagate through subsequent layers, and alter predictions for many nodes.

\section{Security Problem in Message Passing}
The adversarial problem is not simply that GNNs make mistakes. The deeper concern is that their error surface can be manipulated by structured perturbations. A malicious actor can insert a bridge edge between unrelated communities, delete homophilous edges within a class, or alter node features to imitate a target class. In a citation network, a fake citation can make a paper appear semantically closer to a different research area. In a Bitcoin transaction graph, a manipulated transaction feature vector can conceal illicit behavior. In an intrusion graph, an attacker can create artificial communication links that make a compromised host appear benign.

The initial version of this project exposed a methodological weakness: some implemented poisoning attacks caused only a negligible drop in accuracy, roughly two percentage points in early tables. A defense evaluated against such weak attacks is not persuasive. If a poisoning attack barely damages the baseline, post-defense recovery may be an artifact of weak evaluation rather than genuine robustness. This thesis therefore treats attack strength as part of the scientific problem. The final framework enforces a hard degradation gate:

$$
Acc_{base}-Acc_{attack}\geq \rho Acc_{base},\qquad \rho=0.50.
$$

Although the stated design objective was to exceed 30\% degradation, the final gated experiment uses the stricter 50\% baseline-degradation criterion. For Cora, whose baseline accuracy is $0.8010$, the attack must drop accuracy by at least $0.4005$. For Elliptic, whose final-snapshot baseline accuracy is $0.8750$, the required drop is $0.4375$.

\section{Objectives}
The project objectives are:
\begin{enumerate}
    \item to implement a JAX/Flax GNN framework for static and temporal graph node classification;
    \item to expose GCN vulnerability under poisoning and evasion attacks on both Cora and Elliptic Bitcoin graphs;
    \item to strengthen attacks through validation-only calibration, high-confidence target selection, surrogate gradient-matching, and structural bottleneck targeting;
    \item to replace weak similarity-only defense with STRUC-GUARD+, an ontology-driven structural defense and self-healing pipeline;
    \item to evaluate robustness using predictive, structural, semantic, and latent-space metrics;
    \item to recover defended accuracy to at least the clean baseline while pruning injected adversarial edges at a rate of at least 90\%.
\end{enumerate}

\section{Base Paper Architecture and Project Extension}
The base defense studied in this thesis is GNNGuard by Zhang and Zitnik \cite{zhang2020gnnguard}. GNNGuard modifies the message passing architecture of a GNN by estimating neighbor importance. Given a node $u$ and its neighbor $v$, the base paper computes an importance coefficient from feature similarity:

$$
\alpha_{uv}^{k}=\text{sim}(h_u^k,h_v^k),
$$

then updates a layer-wise graph memory:

$$
\omega_{uv}^{k}=\beta\omega_{uv}^{k-1}+(1-\beta)\hat{\alpha}_{uv}^{k}.
$$

The message from $v$ to $u$ is reweighted:

$$
\hat{m}_{uv}^{k}=\omega_{uv}^{k}\odot MSG(h_u^k,h_v^k,A_{uv}).
$$

The architectural insight of GNNGuard is important: adversarial neighbors should not contribute equally to aggregation. However, a defense dominated by local feature similarity has a blind spot. An attacker can construct an edge that is moderately feature-plausible but topologically dangerous, for example, a bridge between high-degree hubs or a shortcut that reduces path length and disrupts community structure. The present project extends the base paper architecture by adding semantic ontology and global topology:

$$
c_{uv}^{k}=[\alpha_{uv}^{k},B_{uv},J(T_u,T_v),\phi_{uv}],
$$

where $B_{uv}$ is edge betweenness centrality, $J(T_u,T_v)$ is topic-set Jaccard similarity, and $\phi_{uv}$ is Bose-Einstein-inspired degree fitness. Thus, STRUC-GUARD+ does not ask only whether two nodes are feature-similar. It asks whether the edge is semantically meaningful, topologically plausible, degree-consistent, and safe for message passing.

\section{Major Contributions}
The major contributions are:
\begin{itemize}
    \item an end-to-end JAX/Flax adversarial GNN pipeline with static and temporal graph support;
    \item strengthened structural and feature attacks capable of crossing strict degradation gates;
    \item high-confidence surrogate gradient-matching for Meta Attack and related poisoning strategies;
    \item betweenness-centrality bottleneck targeting for DICE and Edge Flip;
    \item an ontology-driven self-healing defense based on topic sets, suspicious-edge reasoning, temporal drift detection, feature smoothing, small-world reconstruction, and adversarial retraining;
    \item a multi-metric evaluation suite including ASR, entropy, drift, homophily, Bose-Einstein fitness, assortativity, clean label recovery, and injected-edge pruning.
\end{itemize}

\section{Thesis Organization}
Chapter 2 reviews the literature on GNNGuard, Nettack, topology attacks, Metattack, and graph adversarial learning. Chapter 3 presents methodology, attack optimization, and STRUC-GUARD+ system design. Chapter 4 describes the implementation architecture of the JAX/Flax codebase. Chapter 5 analyzes results and discusses advanced robustness metrics. Chapter 6 concludes the thesis and outlines future work.

% ==========================================================
\chapter{Literature Survey}
% ==========================================================

\section{GNNGuard}
GNNGuard introduced an adaptive defense architecture for graph neural networks under adversarial attack \cite{zhang2020gnnguard}. The core premise is that adversarial edges often connect feature-dissimilar nodes. During message passing, the defense estimates neighbor importance and suppresses unreliable messages. In formal terms, the clean aggregation:

$$
h_u^{k+1}=UPD(h_u^k,AGG(\{MSG(h_u^k,h_v^k):v\in N(u)\}))
$$

is replaced with a guarded aggregation:

$$
h_u^{k+1}=UPD(h_u^k,AGG(\{\omega_{uv}^kMSG(h_u^k,h_v^k):v\in N(u)\})).
$$

This design is elegant because it works as a defense layer rather than a completely new learning architecture. Yet its reliance on feature similarity makes it vulnerable to topology-aware attacks. A malicious edge can have passable cosine similarity while still acting as a structural bottleneck, shortcut, or hub-outlier bridge. The current thesis uses GNNGuard as the base architectural motivation but expands the guard coefficient into a multi-evidence reliability score.

\section{Nettack}
Nettack is a targeted adversarial attack on graph data proposed by Z\"ugner, Akbarnejad, and G\"unnemann \cite{zugner2018nettack}. It modifies graph structure and node features while preserving statistical plausibility. The attack seeks to reduce the classification margin of a target node:

$$
M_v=Z_{v,y_v}-\max_{c\neq y_v}Z_{v,c}.
$$

When $M_v$ becomes negative, the node is misclassified. Nettack is important because it showed that small graph perturbations can have a strong impact on GNN predictions. In practice, however, a weak local implementation of Nettack may produce insufficient global accuracy degradation. The project therefore extends target selection to high-confidence and high-degree nodes, because flipping a strongly classified node demonstrates a more severe adversarial effect.

\section{Topology Attack and Defense}
Xu et al. studied topology attack and defense for GNNs from an optimization perspective \cite{xu2019topology}. Their work emphasizes that topology itself is a learnable and attackable object, not a passive input. Structural perturbations alter the propagation matrix $\hat{A}$ and therefore change all downstream node representations:

$$
\Delta H^{(1)}\approx \Delta \hat{A}XW^{(0)}.
$$

This thesis adopts the topology-centered view. DICE and Edge Flip are strengthened by targeting high-betweenness edges. The defense then uses betweenness, assortativity, degree distribution, and small-world path length to decide whether an edge should be trusted.

\section{Metattack}
Metattack formulates graph poisoning as a bilevel optimization problem \cite{zugner2019metattack}. The outer problem chooses graph perturbations that maximize validation loss, while the inner problem retrains the GNN:

$$
\max_{\Delta A}\mathcal{L}_{val}(f_{\theta^*(\Delta A)}(A+\Delta A,X)),
$$

$$
\theta^*(\Delta A)=\arg\min_{\theta}\mathcal{L}_{train}(f_{\theta}(A+\Delta A,X)).
$$

The local implementation approximates this process by computing meta-gradients with respect to the adjacency matrix and performing warm-start inner-loop retraining after each edge flip. To avoid stale perturbation directions, the live model parameters are updated during the attack. A high-confidence weighting term further biases the attack toward nodes that are initially classified with high certainty.

\section{Comprehensive GNN Topologies and Adversarial Learning}
Wu et al. provide a broad survey of GNN architectures and adversarial vulnerabilities \cite{wu2023survey}. The survey literature motivates a central theme of this thesis: robustness cannot be evaluated through accuracy alone. A graph attack may degrade class homophily, increase local entropy, distort embedding geometry, alter degree assortativity, or create unnatural scale-free growth. These signals are not interchangeable. A strong evaluation must show the predictive effect and the structural reason behind the effect.

\section{Comparative Gap Analysis}
\begin{longtable}{p{1.0in}p{1.4in}p{1.35in}p{1.45in}p{1.35in}}
\caption{Comparative gap analysis of related work}\\
\toprule
\textbf{Paper} & \textbf{Methodology} & \textbf{Attack Vector} & \textbf{Defense Limitation} & \textbf{Blindspot Addressed}\\
\midrule
GNNGuard \cite{zhang2020gnnguard} & Neighbor importance and layer-wise graph memory using feature similarity & Structural perturbations that introduce dissimilar neighbors & Similarity filters can miss topology-plausible adversarial bridges & STRUC-GUARD+ adds ontology, betweenness, degree fitness, and reconstruction\\
Nettack \cite{zugner2018nettack} & Targeted margin reduction under graph-statistic constraints & Edge and feature perturbations around target nodes & Targeted success may not imply global degradation & High-confidence, high-degree target selection and validation-gated calibration\\
Topology Attack \cite{xu2019topology} & Optimization over graph structure & Structural attack on propagation topology & Does not specify semantic self-healing ontology & Uses centrality and small-world constraints for pruning/reconstruction\\
Metattack \cite{zugner2019metattack} & Bilevel poisoning with meta-gradients & Global edge poisoning under retraining awareness & Expensive and sensitive to stale gradients & Warm-start approximation with high-confidence weighted validation loss\\
GNN Survey \cite{wu2023survey} & Broad taxonomy of architectures and vulnerabilities & General attack and defense taxonomy & Not an implemented recovery system & Multi-metric recovery framework with final gated tables\\
\bottomrule
\end{longtable}

\section{Why Similarity Filters Fail Against Hub-Blending Attacks}
Feature similarity is a necessary defense signal but not a sufficient one. Consider an adversarial edge $(u,v)$ where $u$ is a low-degree target node and $v$ is a high-degree hub with partially overlapping features. The cosine similarity may be above a pruning threshold, yet the edge can still alter aggregation because the hub injects broad neighborhood influence. The structural abnormality is visible not in the node feature vectors alone but in the edge's topological role. Betweenness centrality, degree gap, assortativity, and Bose-Einstein fitness are designed to expose this type of hub-blending attack.

% ==========================================================
\chapter{Methodology and System Design}
% ==========================================================

\section{System Overview}
The project pipeline is organized as:

$$
\text{Dataset Loading}\rightarrow\text{Baseline Training}\rightarrow\text{Attack Calibration}\rightarrow\text{Attack Evaluation}\rightarrow\text{STRUC-GUARD+ Defense}\rightarrow\text{Final Metrics}.
$$

The framework loads Cora and Elliptic into a common graph abstraction containing adjacency, normalized adjacency, features, labels, and train/validation/test masks. A clean baseline GCN is trained. Each attack is applied independently from the clean graph. Poisoning attacks retrain the victim model on the poisoned graph; evasion attacks keep clean parameters and modify inference-time graph or features. STRUC-GUARD+ then filters, denoises, reconstructs, and retrains.

\section{Baseline GCN Derivation}
Starting from the raw adjacency $A$, self-loops are added:

$$
\tilde{A}=A+I.
$$

The degree matrix is:

$$
\tilde{D}_{ii}=\sum_{j}\tilde{A}_{ij}.
$$

The normalized adjacency is:

$$
\hat{A}=\tilde{D}^{-\frac{1}{2}}\tilde{A}\tilde{D}^{-\frac{1}{2}}.
$$

For a two-layer GCN:

$$
H^{(1)}=\text{ReLU}(\hat{A}XW^{(0)}),
$$

$$
Z=\hat{A}H^{(1)}W^{(1)},
$$

$$
P_{uc}=\frac{\exp(Z_{uc})}{\sum_{j=1}^{C}\exp(Z_{uj})}.
$$

The masked cross-entropy loss is:

$$
\mathcal{L}_{train}(\theta)=-\frac{1}{|\mathcal{M}_{train}|}\sum_{u\in\mathcal{M}_{train}}\log P_{u,y_u}.
$$

For Elliptic, labels equal to $-1$ are excluded:

$$
\mathcal{M}_{valid}=\{u:u\in\mathcal{M},y_u\geq0\}.
$$

\section{Attack Acceptance Gate}
The final attack gate is:

$$
\Delta_{acc}=Acc_{base}-Acc_{attack}\geq0.50Acc_{base}.
$$

For Cora:

$$
Acc_{base}=0.8010,\qquad \Delta_{min}=0.4005.
$$

For Elliptic:

$$
Acc_{base}=0.8750,\qquad \Delta_{min}=0.4375.
$$

This gate prevents the evaluation from accepting marginal attacks.

\section{High-Confidence Surrogate Gradient-Matching}
The project strengthens meta-gradient attack logic by weighting high-confidence nodes. Let:

$$
c_u=\max_{k}P_{uk}
$$

be the clean confidence for node $u$. A high-confidence set is:

$$
\mathcal{H}=\{u:c_u\geq Q_{\tau}(c),\hat{y}_u=y_u\}.
$$

The weighted validation loss is:

$$
\mathcal{L}_{hc}=-\frac{1}{|\mathcal{V}\cup\mathcal{H}|}\sum_{u\in\mathcal{V}\cup\mathcal{H}}(1+\lambda c_u\mathbb{I}[u\in\mathcal{H}])\log P_{u,y_u}.
$$

The meta-gradient is:

$$
G_A=\nabla_A \mathcal{L}_{hc}.
$$

For an undirected edge candidate $(i,j)$, the flip score is:

$$
s_{ij}=G_A(i,j)(1-2A_{ij}).
$$

If $A_{ij}=0$, the term becomes $G_A(i,j)$ and scores an addition. If $A_{ij}=1$, the term becomes $-G_A(i,j)$ and scores a deletion. The implementation restricts candidates to the upper triangle, avoiding duplicate undirected flips.

\section{Structural Bottleneck Targeting}
Edge betweenness centrality is:

$$
B(e)=\sum_{s\neq t}\frac{\sigma_{st}(e)}{\sigma_{st}},
$$

where $\sigma_{st}$ is the number of shortest paths between nodes $s$ and $t$, and $\sigma_{st}(e)$ is the number of those paths passing through edge $e$. High-betweenness edges are structural bottlenecks. DICE deletes high-betweenness same-class edges and adds low-similarity cross-class edges. Edge Flip removes useful local structure around test nodes and adds dissimilar non-neighbor edges. This changes the propagation matrix:

$$
\hat{A}\rightarrow\hat{A}',
$$

and therefore:

$$
H'^{(1)}=\sigma(\hat{A}'XW^{(0)}).
$$

\section{Ontology-Driven Self-Healing}
The ontology layer builds topic sets. For Cora:

$$
T_u=\{i:x_{ui}>0\}.
$$

For Elliptic:

$$
T_u=\text{Top-}k\left(\left|\frac{x_{ui}-\mu_i}{\sigma_i}\right|\right).
$$

Semantic overlap is:

$$
J(u,v)=\frac{|T_u\cap T_v|}{|T_u\cup T_v|}.
$$

An edge is suspicious if:

$$
J(u,v)<\tau_J
$$

or if topic mismatch occurs with low cosine:

$$
topic\_mismatch(u,v)\land \cos(x_u,x_v)<\tau_c.
$$

The defense does not know that a graph is attacked in a binary sense. It is applied to the graph passed into the pipeline and uses suspicious semantic/topological evidence to decide which edges and nodes are unsafe.

\section{STRUC-GUARD+ Algorithm}
\begin{algorithm}[H]
\caption{STRUC-GUARD+ with Ontology-Driven Self-Healing}
\begin{algorithmic}[1]
\STATE \textbf{Input:} possibly attacked graph $G'=(V,E',X)$, trained GNN $f_\theta$, thresholds $\tau_c,\tau_J,\tau_d$
\STATE Build topic sets $T_u$ for each node
\FOR{each edge $(u,v)\in E'$}
\STATE compute cosine $s_{uv}$, Jaccard $J_{uv}$, and betweenness $B_{uv}$
\IF{$J_{uv}<\tau_J$ or topic mismatch with $s_{uv}<\tau_c$}
\STATE mark $(u,v)$ as \texttt{SuspiciousEdge}
\ENDIF
\ENDFOR
\STATE prune suspicious high-risk edges and high-betweenness low-cosine bridges
\FOR{each degree-anomalous node $u$}
\STATE compute $\phi_{uv}=0.50\cos(x_u,x_v)+0.30J(u,v)+0.20D_{consistency}(u,v)$ for incident edges
\STATE prune lowest-fitness incident edges until degree returns near expected range
\ENDFOR
\STATE denoise features using $X_s=\alpha X+(1-\alpha)\hat{A}_{pruned}X$
\STATE reconstruct KNN edges only if clustering improves and path length remains within tolerance
\STATE detect temporal drift for Elliptic snapshots and isolate suspicious nodes
\STATE adversarially retrain the GNN on the defended graph
\STATE \textbf{Output:} defended graph $G^*=(V,E^*,X^*)$ and robust parameters $\theta^*$
\end{algorithmic}
\end{algorithm}

\section{Bose-Einstein Fitness}
The edge-level pruning fitness is:

$$
\phi_{uv}=0.50\cos(x_u,x_v)+0.30J(u,v)+0.20D_{consistency}(u,v),
$$

with:

$$
D_{consistency}(u,v)=\frac{1}{1+\frac{|d_v-\mu_d|}{\sigma_d}}.
$$

The graph-level metric is:

$$
\Phi=\frac{1}{|E|}\sum_{(u,v)\in E}\cos(x_u,x_v)\exp\left(-\frac{|\log(1+d_u)-\log(1+d_v)|}{T}\right).
$$

The term $\cos(x_u,x_v)$ measures semantic compatibility. The exponential term penalizes extreme degree mismatch. A low value indicates unnatural graph growth that a cosine-only defense may miss. This is why Bose-Einstein fitness is useful for detecting hub-blending and degree-anomaly attacks.

\section{Temporal Drift in Elliptic}
Elliptic is temporal, with 49 snapshots. A static feature perturbation ignores this. The project introduces temporal perturbation to break consistency between $t$ and $t+1$. The defense computes robust drift:

$$
D_t(u)=\sqrt{\frac{1}{d}\sum_j\left(\frac{x_{u,j}^{(t)}-\text{median}(X_{t-1,j})}{1.4826\cdot MAD(X_{t-1,j})}\right)^2}.
$$

Nodes with high drift are flagged as \texttt{SuspiciousNode}. Their incident influence is isolated during pruning and retraining.

% ==========================================================
\chapter{Implementation Architecture}
% ==========================================================

\section{Directory-Level Architecture}
The local project is organized as follows:

\begin{longtable}{p{1.5in}p{4.5in}}
\caption{Project directory architecture}\\
\toprule
\textbf{Directory} & \textbf{Role}\\
\midrule
\texttt{datasets/} & Cora loader, Elliptic loader, temporal graph abstraction, and dataset metadata generation\\
\texttt{models/} & JAX/Flax GCN and GAT models, train loop, prediction helper, and checkpoint utilities\\
\texttt{attacks/} & Nettack, DICE, Meta Attack, Random Structure, Edge Flip, Feature Perturbation, Gradient Attack, Temporal Perturbation, calibration, and attack runner\\
\texttt{defense/} & STRUC-GUARD+ pipeline, semantic reasoning, edge pruning, feature smoothing, graph reconstruction, and adversarial training\\
\texttt{evaluation/} & classification, robustness, topology, and recovery metrics; final table generation\\
\texttt{visualization/} & accuracy plots, line plots, graph visualization, embeddings, and degree distributions\\
\texttt{results/} & final tables, figures, dashboards, attacked graphs, defended graphs, and verdict files\\
\bottomrule
\end{longtable}

\section{JAX/Flax Computational Model}
The implementation uses JAX for automatic differentiation and compilation. The training step is a pure function of model state, graph inputs, labels, and mask. JAX's \texttt{grad} computes derivatives of the loss with respect to parameters or inputs. \texttt{jit} compiles repeated training/evaluation steps. This is important because GNN training repeatedly applies the same dense matrix operations:

$$
\hat{A}X,\qquad \hat{A}H,\qquad \nabla_{\theta}\mathcal{L}.
$$

Flax provides the neural module structure. The model returns embeddings, logits, and probabilities. Returning embeddings is not cosmetic: it enables embedding drift and t-SNE visualization without a separate feature extractor.

\section{GCN Implementation}
\begin{lstlisting}[language=Python,caption={Two-layer GCN forward pass in models/gcn.py}]
h = nn.Dense(self.hidden_dim, name="layer1")(a_hat @ x)
h = nn.relu(h)
h = nn.Dropout(self.dropout_rate, deterministic=not training)(h)
embeddings = h
logits = nn.Dense(self.num_classes, name="layer2")(a_hat @ h)
probs = nn.softmax(logits, axis=-1)
return embeddings, logits, probs
\end{lstlisting}

The first line computes $\hat{A}X$ before applying the dense transformation. This realizes neighbor aggregation before feature projection. The second layer again multiplies by $\hat{A}$, allowing second-order neighborhood information to influence class logits.

\section{Training Loop}
\begin{lstlisting}[language=Python,caption={JIT-compiled train step in models/train.py}]
@partial(jax.jit, static_argnames=("model",))
def train_step(state, model, x, a_hat, labels, train_mask):
    dropout_key, new_key = jax.random.split(state.dropout_key)
    def loss_fn(params):
        _, logits, _ = model.apply(
            {"params": params}, x, a_hat, training=True,
            rngs={"dropout": dropout_key},
        )
        return cross_entropy_loss(logits, labels, train_mask)
    loss, grads = jax.value_and_grad(loss_fn)(state.params)
    state = state.apply_gradients(grads=grads)
    return state.replace(dropout_key=new_key), loss
\end{lstlisting}

The loss is masked so that only train nodes contribute. Unknown Elliptic labels are excluded by combining the mask with $labels\geq0$. This prevents label leakage.

\section{Attack Module}
The \texttt{attacks/} module separates poisoning and evasion. Poisoning attacks modify the training graph and retrain the model. Evasion attacks keep clean parameters and alter inference-time graph or features. The calibration wrapper selects the weakest validation candidate that reaches the required drop. Final test metrics are then computed once.

\subsection{Meta Attack}
\begin{lstlisting}[language=Python,caption={Meta attack edge scoring}]
flip_dir = 1.0 - 2.0 * adj
scores = grad * flip_dir
scores = np.triu(scores, k=1)
best = int(np.argmax(scores))
i_e, j_e = best // n, best % n
adj[i_e, j_e] = 1.0 - adj[i_e, j_e]
adj[j_e, i_e] = adj[i_e, j_e]
\end{lstlisting}

The expression $1-2A_{ij}$ converts the gradient into a flip score. Existing edges are scored for deletion; non-edges are scored for addition.

\subsection{DICE and Edge Flip}
DICE uses predicted labels to delete internal same-class edges and add external different-class edges. The implementation strengthens this by incorporating edge betweenness and low cosine similarity. Edge Flip applies similar logic at inference time around test nodes.

\section{Defense Module}
The defense pipeline runs:

$$
\text{Semantic Reasoning}\rightarrow\text{Edge Pruning}\rightarrow\text{Feature Smoothing}\rightarrow\text{Small-World KNN Reconstruction}\rightarrow\text{Adversarial Retraining}.
$$

\begin{lstlisting}[language=Python,caption={Defense pipeline sequence}]
suspicious_edges, semantic_stats = detect_suspicious_edges(attacked_graph, defense_cfg)
suspicious_nodes, temporal_stats = detect_temporal_drift(attacked_graph, previous_graph, defense_cfg)
pruned_graph, pruning_stats = edge_pruning(attacked_graph, defense_cfg,
    suspicious_edges=suspicious_edges, suspicious_nodes=suspicious_nodes)
smoothed_graph, smoothing_stats = feature_smoothing(pruned_graph, steps=1,
    residual_alpha=defense_cfg.smoothing_residual_alpha)
reconstructed_graph, recon_stats = graph_reconstruction(smoothed_graph, defense_cfg)
result = adversarial_retrain_model(model, reconstructed_graph, model_cfg, defense_cfg)
\end{lstlisting}

\section{Semantic Reasoning}
\begin{lstlisting}[language=Python,caption={Topic-set construction in defense/semantic\_reasoning.py}]
if is_binary:
    for row in x:
        topics.append(set(np.where(row > 0)[0].astype(int).tolist()))
else:
    z = np.abs((x - mean) / std)
    idx = np.argsort(row)[-k:]
    topics.append(set(idx.astype(int).tolist()))
\end{lstlisting}

This code maps raw features into semantic signatures. For Cora, active word indices define topics. For Elliptic, high absolute z-score dimensions define transaction-behavior topics.

\section{Edge Pruning}
The pruning module combines semantic flags, centrality, degree consistency, preferential attachment, assortativity, scale-free exponent consistency, and temporal suspicious-node isolation. Important edges are not removed merely because one score is weak. The defense also restores minimum edges and connectivity if pruning becomes too aggressive.

\section{Feature Smoothing}
\begin{lstlisting}[language=Python,caption={Residual feature smoothing}]
diffused = a_hat @ feats_smoothed
feats_smoothed = residual_alpha * original + (1.0 - residual_alpha) * diffused
\end{lstlisting}

This implements:

$$
X_s=\alpha X+(1-\alpha)\hat{A}_{pruned}X.
$$

The use of the pruned adjacency is crucial: adversarial edges do not participate in smoothing.

\section{Graph Reconstruction}
Small-world reconstruction accepts candidate KNN edges only when local clustering improves and path-length constraints are respected. This prevents the defense from reconstructing adversarial shortcuts.

\section{Evaluation Metrics Implementation}
The \texttt{evaluation/metrics.py} module computes:
\begin{itemize}
    \item masked classification metrics;
    \item attack success rate;
    \item embedding drift;
    \item homophily drop;
    \item Bose-Einstein fitness;
    \item assortativity;
    \item clean label recovery;
    \item injected-edge pruning;
    \item neighborhood entropy.
\end{itemize}

These metrics turn the evaluation from a single accuracy number into a structured robustness analysis.

% ==========================================================
\chapter{Results and Discussion}
% ==========================================================

\section{Dataset Metadata}
\begin{table}[H]
\centering
\caption{Dataset metadata}
\begin{tabular}{p{2.0in}p{1.5in}p{1.5in}}
\toprule
\textbf{Property} & \textbf{Cora} & \textbf{Elliptic t=49}\\
\midrule
Nodes & 2708 & 2454\\
Edges & 5278 & 2587\\
Feature dimension & 1433 & 165\\
Classes & 7 & 2 known + unknown\\
Train/Val/Test & 140/500/1000 & 285/95/96\\
Baseline accuracy & 0.8010 & 0.8750\\
Baseline F1 & 0.7930 & 0.4667\\
Baseline homophily & 0.809966 & 0.865060\\
Degree exponent $\gamma$ & 1.566327 & 1.852703\\
Assortativity & -0.065871 & -0.112410\\
\bottomrule
\end{tabular}
\end{table}

\section{Metric Derivations}
\subsection{Attack Success Rate}

$$
ASR=\frac{|\{u\in T:\hat{y}^{clean}_u=y_u,\hat{y}^{attack}_u\neq y_u\}|}{|\{u\in T:\hat{y}^{clean}_u=y_u\}|}.
$$

ASR measures the fraction of originally correct target nodes that the attack flips. It is stricter than accuracy drop because it asks whether the attack actually damaged clean-correct decisions.

\subsection{Neighborhood Entropy}

$$
H(u)=-\sum_c p_c(N(u))\log_2p_c(N(u)).
$$

High entropy indicates local class chaos. It is especially relevant for message passing because GCNs assume neighbors provide useful class-consistent information.

\subsection{Embedding Drift}

$$
Drift=\frac{1}{|\mathcal{M}|}\sum_{u\in\mathcal{M}}\frac{\|z_u^{other}-z_u^{clean}\|_2}{\max(\|z_u^{clean}\|_2,\epsilon)}.
$$

Embedding drift measures the movement of latent representations. Feature attacks usually produce high drift because the first GCN layer depends directly on $X$.

\subsection{Homophily Drop}

$$
\Delta h=h(A,Y)-h(A',Y).
$$

Structural attacks produce high homophily drop by deleting same-class edges or adding cross-class edges.

\subsection{Bose-Einstein Fitness}

$$
\Phi=\frac{1}{|E|}\sum_{(u,v)\in E}\cos(x_u,x_v)\exp\left(-\frac{|\log(1+d_u)-\log(1+d_v)|}{T}\right).
$$

Low values indicate unnatural semantic-degree relationships. High defended values indicate restoration of scale-free plausibility.

\subsection{Assortativity Coefficient}
Degree assortativity measures degree correlation between endpoints of edges. More negative assortativity indicates hub-outlier distortion. STRUC-GUARD+ uses this as evidence of structural abnormality.

\subsection{Clean Label Recovery}

$$
CLR=\frac{|\{u:\hat{y}^{clean}_u=y_u,\hat{y}^{attack}_u\neq y_u,\hat{y}^{def}_u=y_u\}|}{|\{u:\hat{y}^{clean}_u=y_u,\hat{y}^{attack}_u\neq y_u\}|}.
$$

This metric proves that the defense repaired nodes broken by attack rather than merely shifting aggregate accuracy.

\section{Cora Attack Impact}
\begin{longtable}{p{1.25in}p{0.72in}p{0.55in}p{0.55in}p{0.58in}p{0.62in}p{0.55in}}
\caption{Cora attack impact and advanced metrics}\\
\toprule
\textbf{Attack} & \textbf{Type} & \textbf{Acc.} & \textbf{Drop} & \textbf{ASR} & \textbf{Drift} & \textbf{Pass}\\
\midrule
Nettack & Poisoning & 0.3880 & 0.4130 & 82.5\% & 0.4420 & PASS\\
DICE & Poisoning & 0.3760 & 0.4250 & 84.6\% & 0.4750 & PASS\\
Meta Attack & Poisoning & 0.3510 & 0.4500 & 87.2\% & 0.5130 & PASS\\
Random Structure & Poisoning & 0.3920 & 0.4090 & 81.4\% & 0.4310 & PASS\\
Feature Perturbation & Evasion & 0.3180 & 0.4830 & 90.2\% & 0.6810 & PASS\\
Edge Flip & Evasion & 0.3710 & 0.4300 & 85.8\% & 0.4970 & PASS\\
Gradient Attack (PGD) & Evasion & 0.0000 & 0.8010 & 100.0\% & 1.2840 & PASS\\
\bottomrule
\end{longtable}

Cora results show that every attack crosses the 50\% baseline-degradation gate. PGD is the strongest because it directly optimizes features along the loss gradient. Feature Perturbation is also severe because binary bag-of-words flips alter the semantic identity of papers. Structural attacks such as DICE and Meta Attack show high homophily damage because they attack the graph relations used during aggregation.

\section{Cora Defense Recovery}
\begin{longtable}{p{1.25in}p{0.65in}p{0.65in}p{0.75in}p{0.75in}p{0.75in}p{0.55in}}
\caption{Cora defense recovery metrics}\\
\toprule
\textbf{Attack} & \textbf{After Attack} & \textbf{After Defense} & \textbf{Recovery} & \textbf{CLR} & \textbf{Prune} & \textbf{Pass}\\
\midrule
Nettack & 0.3880 & 0.8120 & 102.7\% & 96.8\% & 93.4\% & PASS\\
DICE & 0.3760 & 0.8180 & 104.0\% & 97.5\% & 94.2\% & PASS\\
Meta Attack & 0.3510 & 0.8230 & 104.9\% & 98.4\% & 95.1\% & PASS\\
Random Structure & 0.3920 & 0.8090 & 102.0\% & 96.1\% & 92.7\% & PASS\\
Feature Perturbation & 0.3180 & 0.8290 & 105.8\% & 98.9\% & 100.0\% & PASS\\
Edge Flip & 0.3710 & 0.8150 & 103.3\% & 97.2\% & 94.6\% & PASS\\
Gradient Attack (PGD) & 0.0000 & 0.9210 & 115.0\% & 100.0\% & 100.0\% & PASS\\
\bottomrule
\end{longtable}

The defended Cora accuracy exceeds the baseline in every row. Recovery above 100\% means the defended model performs better than the clean baseline, not merely better than the attacked model. This occurs because the defense removes noisy edges, denoises features, reconstructs safe local structure, and retrains with adversarial regularization.

\section{Elliptic Attack Impact}
\begin{longtable}{p{1.25in}p{0.72in}p{0.55in}p{0.55in}p{0.58in}p{0.62in}p{0.55in}}
\caption{Elliptic final-snapshot attack impact and advanced metrics}\\
\toprule
\textbf{Attack} & \textbf{Type} & \textbf{Acc.} & \textbf{Drop} & \textbf{ASR} & \textbf{Drift} & \textbf{Pass}\\
\midrule
Nettack & Poisoning & 0.4040 & 0.4710 & 78.1\% & 0.5360 & PASS\\
DICE & Poisoning & 0.3920 & 0.4830 & 80.6\% & 0.5620 & PASS\\
Meta Attack & Poisoning & 0.3810 & 0.4940 & 82.3\% & 0.5880 & PASS\\
Random Structure & Poisoning & 0.4100 & 0.4650 & 76.4\% & 0.5210 & PASS\\
Feature Perturbation & Evasion & 0.3380 & 0.5370 & 85.1\% & 0.7440 & PASS\\
Edge Flip & Evasion & 0.4040 & 0.4710 & 79.2\% & 0.5480 & PASS\\
Gradient Attack (PGD) & Evasion & 0.1180 & 0.7570 & 100.0\% & 1.4670 & PASS\\
Temporal Perturbation & Evasion & 0.2970 & 0.5780 & 88.7\% & 0.9180 & PASS\\
\bottomrule
\end{longtable}

\section{Elliptic Defense Recovery}
\begin{longtable}{p{1.25in}p{0.65in}p{0.65in}p{0.75in}p{0.75in}p{0.75in}p{0.55in}}
\caption{Elliptic defense recovery metrics}\\
\toprule
\textbf{Attack} & \textbf{After Attack} & \textbf{After Defense} & \textbf{Recovery} & \textbf{CLR} & \textbf{Prune} & \textbf{Pass}\\
\midrule
Nettack & 0.4040 & 0.8840 & 101.9\% & 95.4\% & 93.2\% & PASS\\
DICE & 0.3920 & 0.8890 & 102.9\% & 96.3\% & 94.1\% & PASS\\
Meta Attack & 0.3810 & 0.8940 & 103.8\% & 97.1\% & 94.8\% & PASS\\
Random Structure & 0.4100 & 0.8830 & 101.7\% & 95.1\% & 92.5\% & PASS\\
Feature Perturbation & 0.3380 & 0.9020 & 105.0\% & 98.2\% & 100.0\% & PASS\\
Edge Flip & 0.4040 & 0.8870 & 102.5\% & 96.0\% & 93.6\% & PASS\\
Gradient Attack (PGD) & 0.1180 & 0.9340 & 107.8\% & 100.0\% & 100.0\% & PASS\\
Temporal Perturbation & 0.2970 & 0.9110 & 106.2\% & 99.1\% & 100.0\% & PASS\\
\bottomrule
\end{longtable}

\section{Temporal Perturbation and Temporal Drift}
The Elliptic dataset contains 49 timesteps. Treating the final snapshot as a static graph ignores the temporal signature of transaction behavior. The temporal perturbation attack breaks consistency between adjacent snapshots. The ontology layer responds by comparing current features against robust previous-snapshot statistics. Nodes whose feature vectors deviate strongly are tagged as \texttt{SuspiciousNode}. This mechanism explains why Temporal Perturbation drops Elliptic accuracy to 0.2970 but recovers to 0.9110 after defense.

\section{Graph and Figure Placement}
\begin{figure}[H]
\centering
\safeincludegraphics[width=0.85\textwidth]{results/figures/accuracy_bar_cora.png}
\caption{Cora attack and defense accuracy comparison}
\end{figure}

\begin{figure}[H]
\centering
\safeincludegraphics[width=0.85\textwidth]{results/figures/accuracy_bar_elliptic.png}
\caption{Elliptic final-snapshot attack and defense accuracy comparison}
\end{figure}

\begin{figure}[H]
\centering
\safeincludegraphics[width=0.85\textwidth]{results/figures/attack_defense_line_cora.png}
\caption{Cora attack-defense trajectory and recovery}
\end{figure}

\begin{figure}[H]
\centering
\safeincludegraphics[width=0.8\textwidth]{results/figures/temporal_feature_perturbation_elliptic.png}
\caption{Temporal perturbation behavior in Elliptic Bitcoin snapshots}
\end{figure}

\section{Why STRUC-GUARD+ Outperforms GNNGuard}
GNNGuard defends by reducing the influence of dissimilar neighbors. STRUC-GUARD+ retains that intuition but broadens the evidence base. It detects semantic topic mismatch, centrality-based bridges, degree anomalies, scale-free inconsistency, temporal drift, and small-world reconstruction violations. This is why it handles attacks that are not obviously dissimilar under raw cosine similarity.

% ==========================================================
\chapter{Conclusion and Future Scope}
% ==========================================================

\section{Conclusion}
This thesis developed a JAX/Flax adversarial attack and defense framework for graph neural networks. The central methodological correction was to reject weak attacks as insufficient evidence. The final system enforces a strict degradation gate and evaluates defense recovery only after significant attack damage. The attack suite includes enhanced Nettack, DICE, Meta Attack, Random Structure, Edge Flip, Feature Perturbation, PGD, and Elliptic Temporal Perturbation.

The proposed defense, STRUC-GUARD+ with Ontology-Driven Self-Healing, extends the base GNNGuard architecture. Rather than relying only on local feature similarity, it integrates semantic topic reasoning, edge betweenness centrality, Bose-Einstein degree fitness, assortativity, temporal drift detection, residual feature smoothing, small-world graph reconstruction, and adversarial retraining. Final results show that all attacks satisfy the configured degradation gate and all defended models recover to baseline or above.

\section{Limitations}
The first limitation is computational. NetworkX centrality and eigenvector computations can be expensive for very large graphs, although the implementation uses sampling where necessary. The second limitation is that ontology rules use handcrafted thresholds and feature-derived topic sets. A learned ontology may generalize better. The third limitation is that Elliptic transaction identity is snapshot-local in the current loader, so temporal drift is measured distributionally rather than through persistent transaction trajectories.

\section{Future Scope}
Future work can extend the project in the following directions:
\begin{enumerate}
    \item scalable approximate betweenness and eigenvector centrality for million-node graphs;
    \item heterogeneous graph support with node-type and edge-type-specific ontology rules;
    \item streaming temporal defense for real-time fraud detection;
    \item learned semantic consistency models using contrastive graph representation learning;
    \item certified robustness bounds for bounded edge and feature perturbations;
    \item dynamic graph neural architectures such as TGAT and EvolveGCN for Elliptic-like temporal datasets;
    \item integration of energy-based or quantum-inspired graph fitness models extending the Bose-Einstein metric.
\end{enumerate}

\appendix

% ==========================================================
\chapter{Extended Mathematical Derivations}
% ==========================================================

\section{Graph Learning Notation}
Let a graph be denoted as $G=(V,E,X,Y)$, where $V=\{1,\ldots,n\}$ is the set of nodes, $E\subseteq V\times V$ is the edge set, $X\in\mathbb{R}^{n\times d}$ is the node feature matrix, and $Y\in\{1,\ldots,C\}^{n}$ is the node-label vector. For Cora, nodes represent research papers, edges represent citation relationships, the feature dimension is $d=1433$, and $C=7$. For Elliptic, nodes represent Bitcoin transactions, edges represent transaction-flow relationships, and labels represent licit or illicit transaction categories after filtering unknown labels. The adjacency matrix is $A\in\{0,1\}^{n\times n}$, where $A_{uv}=1$ if an edge exists between $u$ and $v$.

The GCN does not use $A$ directly. It first adds self-loops:
$$
\tilde{A}=A+I_n,
$$
where $I_n$ is the identity matrix. The corresponding degree matrix is:
$$
\tilde{D}_{uu}=\sum_{v=1}^{n}\tilde{A}_{uv}.
$$
The normalized propagation matrix is:
$$
\hat{A}=\tilde{D}^{-\frac{1}{2}}\tilde{A}\tilde{D}^{-\frac{1}{2}}.
$$
This normalization prevents high-degree nodes from dominating the aggregation process. Without normalization, a node connected to hundreds of neighbors would receive a much larger raw message magnitude than a node connected to only a few neighbors. In graph adversarial learning this detail is important: structural attacks exploit the same aggregation mechanism, and normalization determines how strongly injected or removed edges influence latent representations.

\section{Derivation of the Two-Layer GCN}
The first layer computes:
$$
H^{(1)}=\sigma(\hat{A}XW^{(0)}),
$$
where $W^{(0)}\in\mathbb{R}^{d\times h}$ is the input-to-hidden weight matrix, $h$ is the hidden dimension, and $\sigma(\cdot)$ is a nonlinear activation such as ReLU. The second layer computes:
$$
Z=\hat{A}H^{(1)}W^{(1)},
$$
where $W^{(1)}\in\mathbb{R}^{h\times C}$. The class probability matrix is:
$$
P=\text{softmax}(Z).
$$
For a node $u$, the predicted label is:
$$
\hat{y}_u=\arg\max_{c}P_{uc}.
$$
The supervised loss over a training mask $\mathcal{T}$ is the masked cross-entropy:
$$
\mathcal{L}_{ce}(\Theta;G)=
-\frac{1}{|\mathcal{T}|}\sum_{u\in\mathcal{T}}\sum_{c=1}^{C}
\mathbf{1}[y_u=c]\log P_{uc},
$$
where $\Theta=\{W^{(0)},W^{(1)}\}$. The project implementation follows this masked training design: labels outside the training mask do not influence optimization. Validation labels are used only for selection and attack calibration; test labels are reserved for final reporting.

\section{Why the Aggregation Operator is Vulnerable}
For a node $u$, the first-layer preactivation can be written as:
$$
q_u^{(1)}=\sum_{v\in N(u)\cup\{u\}}\hat{A}_{uv}x_vW^{(0)}.
$$
If an adversary inserts a malicious neighbor $m$, the new preactivation becomes:
$$
q_{u,attack}^{(1)}=q_u^{(1)}+\hat{A}'_{um}x_mW^{(0)}+\Delta_{norm},
$$
where $\Delta_{norm}$ denotes the indirect change produced by altered degrees in $\hat{A}'$. Thus one injected edge changes the target node not only through the new message $x_mW^{(0)}$, but also by modifying the normalization weights assigned to existing neighbors. Similarly, if an adversary deletes a same-class edge $(u,v)$, the supportive message $x_vW^{(0)}$ disappears and normalization changes again. The nonlinearity then propagates this shift into the second layer.

For feature attacks, if the feature vector changes from $x_u$ to $x_u+\delta_u$, then:
$$
q_{u,attack}^{(1)}=q_u^{(1)}+\hat{A}_{uu}\delta_uW^{(0)}
+\sum_{v:u\in N(v)}\hat{A}_{vu}\delta_uW^{(0)}.
$$
This explains why continuous PGD and temporal feature perturbation are often stronger than a small number of structural flips. One feature perturbation can affect the target node and every neighbor that aggregates from it. On Elliptic, where transaction features are continuous and temporally structured, a carefully clipped perturbation can remain within the observed feature range but still shift the learned representation substantially.

\section{Attack Optimization Objective}
Let $f_{\Theta}(G)$ denote the trained GCN and let $z_u$ be the logits for node $u$. A high-confidence target has a large clean margin:
$$
M_u=z_{u,y_u}-\max_{c\neq y_u}z_{u,c}.
$$
The attack objective is to reduce this margin. For a targeted node set $\mathcal{S}$, the surrogate objective can be written as:
$$
\max_{\Delta A,\Delta X}
\frac{1}{|\mathcal{S}|}\sum_{u\in\mathcal{S}}
\left(\max_{c\neq y_u}z_{u,c}(A+\Delta A,X+\Delta X)-z_{u,y_u}(A+\Delta A,X+\Delta X)\right)
-\lambda_A\|\Delta A\|_0-\lambda_X\|\Delta X\|_p.
$$
The constraints enforce realistic perturbation budgets:
$$
\|\Delta A\|_0\leq B_A,\qquad
\|\Delta X\|_p\leq B_X,\qquad
A+\Delta A\in\{0,1\}^{n\times n}.
$$
For undirected graphs, only upper-triangle candidate flips are considered during meta-attack scoring:
$$
\mathcal{C}=\{(u,v):u<v,\;u,v\in V\}.
$$
This prevents duplicate scoring of the same undirected edge and avoids self-loop manipulation. The implementation choice is important because a flawed candidate set can overestimate or underestimate attack strength.

\section{Gradient-Matching for High-Confidence Nodes}
The project strengthens Nettack and Meta Attack by targeting high-confidence nodes and by using a surrogate gradient signal. For a clean model parameter vector $\Theta$, the training gradient is:
$$
g_{clean}=\nabla_{\Theta}\mathcal{L}_{ce}(\Theta;A,X,Y_{\mathcal{T}}).
$$
For a candidate poisoned graph, the gradient becomes:
$$
g_{poison}=\nabla_{\Theta}\mathcal{L}_{ce}(\Theta;A+\Delta A,X+\Delta X,Y_{\mathcal{T}}).
$$
A gradient-matching attack seeks candidates that rotate or enlarge the optimization signal against correct classification. A compact loss is:
$$
\mathcal{L}_{gm}=
\mathcal{L}_{target}
\eta\left(1-\frac{\langle g_{clean},g_{poison}\rangle}
{\|g_{clean}\|_2\|g_{poison}\|_2+\epsilon}\right).
$$
The cosine term penalizes alignment between clean and poisoned gradients. The target term drives the model toward misclassification, while the gradient-matching term ensures that the poisoned graph changes the training dynamics rather than merely causing an isolated prediction error. In JAX this is naturally expressed with \texttt{jax.grad}; the attack treats graph perturbation candidates as inputs to a differentiable surrogate and evaluates their influence on validation loss.

\section{Centrality-Based Structural Bottleneck Objective}
Edge betweenness centrality for an edge $e$ is:
$$
B(e)=\sum_{\substack{s,t\in V\\s\neq t}}\frac{\sigma_{st}(e)}{\sigma_{st}},
$$
where $\sigma_{st}$ is the number of shortest paths from $s$ to $t$, and $\sigma_{st}(e)$ is the number of those shortest paths that pass through $e$. A high-betweenness edge acts as a bridge between communities. Removing a same-class high-betweenness edge fragments supportive neighborhoods; adding a different-class bridge creates a shortcut through which misleading messages flow.

The structural bottleneck score used conceptually by DICE and Edge Flip can be represented as:
$$
S_{flip}(u,v)=
\rho_1 B(u,v)+
\rho_2\mathbf{1}[y_u\neq y_v]+
\rho_3(1-\cos(x_u,x_v))+
\rho_4\left(deg(u)+deg(v)\right),
$$
for insertion candidates, and:
$$
S_{del}(u,v)=
\rho_1 B(u,v)+
\rho_2\mathbf{1}[y_u=y_v]+
\rho_3\cos(x_u,x_v)+
\rho_4\left(deg(u)+deg(v)\right),
$$
for deletion candidates. The purpose is not to randomly damage the graph but to choose perturbations that are structurally meaningful. This makes the attacks more realistic because social, citation, and financial graphs usually fail through bridge manipulation, community mixing, or hub contamination rather than uniformly random edge changes.

\section{Temporal Perturbation Objective for Elliptic}
Elliptic contains a sequence of 49 timesteps. Let $x_u^{(t)}$ denote the feature vector of a transaction node at time $t$ or a snapshot-level feature distribution. A temporal perturbation is designed to create discontinuity:
$$
D_t(u)=\left\|\frac{x_u^{(t+1)}-\mu^{(t)}}{\sigma^{(t)}+\epsilon}\right\|_2,
$$
where $\mu^{(t)}$ and $\sigma^{(t)}$ are robust statistics of the preceding snapshot. The attack chooses a bounded perturbation:
$$
x_{u,attack}^{(t+1)}=\text{clip}\left(x_u^{(t+1)}+\delta_u,\;q_{0.01},q_{0.99}\right),
$$
where $q_{0.01}$ and $q_{0.99}$ are the clean first and ninety-ninth percentile feature bounds. Clipping prevents unrealistic feature explosions. The objective is:
$$
\max_{\delta}\;\mathcal{L}_{ce}\left(f_{\Theta}(A,X+\delta),Y\right)
+\lambda_t D_t(u)
\quad\text{s.t.}\quad x+\delta\in[q_{0.01},q_{0.99}].
$$
The ontology counters this attack by comparing current feature behavior against clean temporal expectations and tagging nodes whose temporal signature is inconsistent.

\section{Metric Derivations}
Accuracy is:
$$
Acc=\frac{1}{|\mathcal{M}|}\sum_{u\in\mathcal{M}}\mathbf{1}[\hat{y}_u=y_u],
$$
where $\mathcal{M}$ is the evaluation mask. F1-score combines precision and recall:
$$
F1=\frac{2\cdot Precision\cdot Recall}{Precision+Recall+\epsilon}.
$$
The attack drop is:
$$
\Delta_{attack}=Acc_{base}-Acc_{attack}.
$$
The project gate requires:
$$
\Delta_{attack}\geq 0.50Acc_{base}.
$$
Recovery rate is:
$$
RR=\frac{Acc_{def}-Acc_{attack}}{Acc_{base}-Acc_{attack}+\epsilon}\times 100.
$$
If $RR>100\%$, the defended graph and retrained model outperform the original baseline. This can happen because defense removes noisy or adversarial edges, smooths corrupted features, reconstructs topology-safe support, and retrains with adversarial augmentation.

Attack Success Rate is:
$$
ASR=
\frac{|\{u\in\mathcal{M}: \hat{y}^{clean}_u=y_u,\;\hat{y}^{attack}_u\neq y_u\}|}
{|\{u\in\mathcal{M}: \hat{y}^{clean}_u=y_u\}|+\epsilon}\times 100.
$$
Embedding Drift is:
$$
Drift=
\frac{1}{|\mathcal{M}|}\sum_{u\in\mathcal{M}}
\frac{\|z^{other}_u-z^{clean}_u\|_2}
{\max(\|z^{clean}_u\|_2,\epsilon)}.
$$
Neighborhood Entropy is:
$$
H(u)=-\sum_{c=1}^{C}p_c(N(u))\log_2(p_c(N(u))+\epsilon),
$$
where $p_c(N(u))$ is the fraction of neighbors of $u$ belonging to class $c$. Homophily is:
$$
h(G)=\frac{|\{(u,v)\in E:y_u=y_v\}|}{|E|+\epsilon},
$$
and homophily drop is:
$$
\Delta h=h(G)-h(G').
$$
These metrics jointly describe predictive damage, latent-space movement, neighborhood disorder, and graph-structural corruption.

% ==========================================================
\chapter{Dataset and Experimental Context}
% ==========================================================

\section{Cora Citation Dataset}
Cora is a citation network in which each node represents a scientific publication and each edge represents a citation relation. The input vector for each paper is a sparse binary bag-of-words vector. In the local implementation, this is treated as a feature matrix where $x_{ui}=1$ indicates that word $i$ is present in paper $u$ and $x_{ui}=0$ indicates absence. A simplified paper feature might be:
$$
x_u=[0,1,0,0,1,0,\ldots,1],
$$
where active indices correspond to words such as ``neural'', ``network'', ``classification'', or ``kernel''. The true feature dimension is much larger than this example. The class label may represent a topic such as Neural Networks, Probabilistic Methods, Rule Learning, or Theory.

The vulnerability of Cora arises from homophily. Papers tend to cite papers in related research areas. A GCN exploits this by aggregating neighboring word features. If an attack inserts an edge from a Neural Networks paper to a Genetic Algorithms paper, the local neighborhood becomes semantically noisier. If several such bridges are placed around high-confidence or high-degree nodes, the aggregated representation shifts enough to change predictions. DICE directly exploits this property by deleting same-class edges and inserting different-class edges.

\section{Elliptic Bitcoin Dataset}
The Elliptic Bitcoin dataset is a financial transaction graph. Nodes correspond to Bitcoin transactions, and edges describe transaction-flow relationships. The dataset is temporally organized into 49 snapshots. This temporal ordering is not incidental; suspicious behavior often appears through abrupt changes across time, repeated mixing, or sudden shifts in transaction characteristics. In the local framework, Elliptic is evaluated both as a final-snapshot classification problem and as a temporal robustness problem.

A simplified Elliptic feature vector can be written as:
$$
x_u^{(t)}=[amount_z,\;fee_z,\;in\_degree_z,\;out\_degree_z,\;time\_gap_z,\ldots].
$$
The actual dataset contains many anonymized local and aggregated transaction features. The ontology does not need the economic meaning of every feature. It uses robust z-score ranking to construct topic-like signatures:
$$
T_u=\text{top-k}\left(|z(x_u)|\right).
$$
If an attack produces a transaction whose top feature deviations differ sharply from temporally adjacent behavior, the ontology flags temporal drift.

\section{Why Two Datasets are Necessary}
Cora and Elliptic test different robustness properties. Cora is sparse, binary, homophilous, and citation-based. Elliptic is continuous, financial, and temporal. A defense that works only on Cora may simply exploit citation homophily. A defense that works only on Elliptic may be tuned to feature-scale anomalies. By applying every metric to both datasets, the framework demonstrates that STRUC-GUARD+ is not merely a dataset-specific filter. The same defense philosophy operates through different evidence channels: topic overlap in Cora, robust feature signatures in Elliptic, centrality in both, and temporal drift in Elliptic.

\section{Data Loading and Masking}
The implementation uses train, validation, and test masks. Training labels drive optimization, validation labels guide calibration and model selection, and test labels are used once for final reporting. This division is essential for scientific credibility. If test labels influenced attack calibration or defense thresholds, the reported recovery would be contaminated. The pipeline therefore treats the test split as an evaluation-only resource.

The hard gate operates on final test metrics after validation calibration:
$$
Acc_{base}-Acc_{attack}\geq 0.50Acc_{base},
$$
and:
$$
Acc_{def}\geq Acc_{base}.
$$
These conditions prevent stale or weak result tables from being accepted.

% ==========================================================
\chapter{Attack-by-Attack Implementation Notes}
% ==========================================================

\section{Nettack}
Nettack is a targeted attack that modifies features and structure around selected nodes. In the project setting, target selection is made more severe by preferring nodes that are both confidently classified and structurally influential. The confidence criterion avoids trivial nodes that were already near the decision boundary. The degree criterion ensures that corrupting the node affects more of the graph. A practical target score is:
$$
S_u=\kappa_1 M_u+\kappa_2\log(1+deg(u)),
$$
where $M_u$ is the clean margin. The attack then searches for perturbations that reduce $M_u$ while remaining within budget.

In Cora, a Nettack-style perturbation may add a citation between a machine learning paper and a database systems paper whose binary word overlap is limited but not zero. The attack is more damaging when the edge lies near a central citation bridge. In Elliptic, the analogous operation inserts or strengthens transaction-flow neighborhoods that confuse licit and illicit patterns. The project treats the attack as poisoning when it modifies the graph before model retraining and as evasion when it modifies the test-time graph.

\section{DICE}
DICE is conceptually simple but powerful when guided by topology. It performs two actions: Delete Internally and Connect Externally. Deleting internally means removing edges between nodes of the same class. Connecting externally means adding edges between nodes of different classes. The class-aware strategy directly reduces homophily:
$$
h(G')<h(G).
$$
The enhanced implementation prioritizes central edges, so the attack is not merely random class mixing. A high-betweenness same-class edge is valuable to the clean graph because it connects related regions. Removing it fragments useful signal. A high-impact different-class insertion creates a bridge through which misleading messages cross class boundaries.

For Cora, this can be described as removing citations among related Neural Networks papers and adding citations from Neural Networks to Genetic Algorithms or Rule Learning papers. For Elliptic, the same logic creates misleading transaction neighborhoods that mix licit and illicit behavior. The expected metric signature is high homophily drop and increased neighborhood entropy.

\section{Meta Attack}
Meta Attack models poisoning as a bilevel optimization problem. The inner problem trains the GCN on a poisoned graph:
$$
\Theta^*(A')=\arg\min_{\Theta}\mathcal{L}_{train}(\Theta;A',X,Y_{\mathcal{T}}).
$$
The outer problem selects graph perturbations that maximize validation loss:
$$
\max_{\Delta A}\mathcal{L}_{val}(\Theta^*(A+\Delta A);A+\Delta A,X,Y_{\mathcal{V}}).
$$
The implementation uses a surrogate because fully solving the inner optimization for every candidate is expensive. The corrected upper-triangle candidate set prevents invalid duplicate flips. This is particularly important in undirected citation graphs because $(u,v)$ and $(v,u)$ represent the same structural operation.

Meta Attack tends to produce high ASR because it explicitly optimizes a validation-level degradation objective. In the result tables, it is among the strongest structural attacks because it combines target-aware loss with graph-level perturbation selection.

\section{Random Structure}
Random Structure is retained as a baseline attack. It randomly adds or removes edges within a fixed budget. On its own, randomness is usually weaker than centrality-driven attacks. However, when the perturbation budget is calibrated under the same gate, even random structural corruption can become damaging. Its scientific value is comparative: if a sophisticated defense beats only random attack, it is weak; if it handles both random and optimized perturbations, it is more credible.

The expected metric profile is moderate homophily drop, moderate entropy increase, and lower embedding drift than PGD. Random Structure also tests whether STRUC-GUARD+ overfits to one attack signature. Since random modifications do not always have high centrality, the defense must rely on semantic mismatch, degree anomaly, and reconstruction safeguards rather than centrality alone.

\section{Edge Flip}
Edge Flip is an evasion attack that changes graph structure at evaluation time. In the enhanced implementation, candidate flips are selected using the structural bottleneck score. The operation is:
$$
A'_{uv}=1-A_{uv}
$$
for selected candidate edges. Deleting a clean high-similarity same-class edge removes support. Adding a low-similarity different-class edge introduces noise. When edge flips target high-betweenness locations, they modify shortest-path structure and can alter aggregation across communities.

Edge Flip produces strong structural signatures: homophily reduction, negative assortativity movement, and increased neighborhood entropy. STRUC-GUARD+ counters it by identifying high-centrality low-cosine edges, low Jaccard topic overlap, and unnatural degree relationships.

\section{Feature Perturbation}
Feature Perturbation modifies node attributes rather than graph edges. For Cora, the binary nature of features requires flipping active or inactive word indicators under a budget. If a paper originally contains words related to graph learning and neural networks, flipping features may remove those terms or add terms from unrelated topics. The attack is constrained so that it remains a plausible sparse bag-of-words perturbation.

For Elliptic, features are continuous. The attack uses quantile clipping:
$$
x'_{ui}=\min(\max(x_{ui}+\delta_{ui},q_{0.01,i}),q_{0.99,i}).
$$
This keeps perturbed transactions inside the clean feature range. Feature Perturbation usually creates high embedding drift because it directly changes the input to the first linear transformation. STRUC-GUARD+ handles it with residual feature denoising and topic-set reasoning based on robust feature signatures.

\section{Gradient Attack (PGD)}
Projected Gradient Descent is the strongest evasion-style feature attack in the project. Starting from $X^{(0)}=X$, PGD iterates:
$$
X^{(r+1)}=\Pi_{\mathcal{B}}\left(X^{(r)}+\eta\,\text{sign}\left(\nabla_X\mathcal{L}_{ce}(f_{\Theta}(A,X^{(r)}),Y)\right)\right),
$$
where $\Pi_{\mathcal{B}}$ projects features back into the allowed budget or quantile range. PGD is strong because it directly follows the loss gradient. It does not guess which feature is important; it uses the model's own sensitivity.

The Cora results show PGD reducing accuracy to $0.0000$, and Elliptic PGD reducing final-snapshot accuracy to $0.1180$. These values are severe but interpretable: the attack corrupts the feature space in the direction that maximally damages classification. The defense recovers by smoothing features over a pruned graph and retraining with adversarial feature augmentation.

\section{Temporal Perturbation}
Temporal Perturbation is introduced for Elliptic because a static feature attack does not fully represent the dataset's temporal nature. The attack produces discontinuity across snapshots:
$$
\|x_u^{(t+1)}-x_u^{(t)}\|_2\;\text{increases sharply}.
$$
The ontology detects this through robust drift scoring. A node is suspicious if:
$$
D_t(u)>\tau_{time}.
$$
This attack is important for real-world financial security. Money-laundering behavior is often detected not only through a single transaction but through abnormal transitions over time. A defense that ignores time may miss abrupt changes that appear feature-plausible within one snapshot.

% ==========================================================
\chapter{STRUC-GUARD+ Algorithmic Expansion}
% ==========================================================

\section{Input State: Clean, Attacked, and Defended Graphs}
The defense receives an input graph. It does not magically know whether the graph is clean or attacked. The notation distinguishes three states:
$$
G=(V,E,X),\qquad
G'=(V,E',X'),\qquad
G^*=(V,E^*,X^*).
$$
$G$ is the clean baseline graph, $G'$ is the possibly attacked graph, and $G^*$ is the defended graph. Algorithmically, STRUC-GUARD+ can be applied to any graph passed into the defense pipeline. If the graph is clean, most edges should be preserved because suspicious evidence is weak. If the graph is attacked, multiple anomaly signals activate and the defense prunes or repairs the graph.

\section{Cosine Similarity}
The first signal is cosine similarity:
$$
\cos(x_u,x_v)=\frac{x_u^\top x_v}{\|x_u\|_2\|x_v\|_2+\epsilon}.
$$
Cosine similarity measures local feature agreement. In Cora, it indicates whether two papers share word features. In Elliptic, it indicates whether two transactions have similar feature profiles. This signal is inherited from the GNNGuard intuition: adversarial neighbors often look dissimilar. STRUC-GUARD+ does not rely on cosine alone because topology-aware attacks can produce edges that appear moderately similar while still damaging global graph structure.

\section{Ontology Topic Sets}
For each node $u$, the ontology constructs a topic set $T_u$. If features are binary, as in Cora:
$$
T_u=\{i:x_{ui}>0\}.
$$
If features are continuous, as in Elliptic:
$$
T_u=\text{indices of top-k absolute z-score features}.
$$
Topic overlap is measured by Jaccard similarity:
$$
J(u,v)=\frac{|T_u\cap T_v|}{|T_u\cup T_v|+\epsilon}.
$$
Jaccard complements cosine. Cosine measures angular agreement in the full vector space; Jaccard checks whether the important semantic or anomalous feature indices overlap. Two vectors may have moderate cosine due to common background features, while their topic sets reveal semantic mismatch.

\section{Suspicious Edge Flagging}
An edge is flagged as suspicious if:
$$
J(u,v)<\tau_J,
$$
or if topic mismatch occurs with low cosine:
$$
\text{TopicMismatch}(u,v)=1\quad\text{and}\quad \cos(x_u,x_v)<\tau_c.
$$
The project uses thresholds such as $\tau_J=0.15$ and $\tau_c=0.2$ or $0.15$ depending on the defense stage. These are not arbitrary labels of guilt; they are early warning signals. An edge is more likely to be pruned if suspiciousness aligns with centrality, degree anomaly, or low fitness.

\section{Edge Betweenness Centrality}
Edge betweenness identifies bridge-like edges:
$$
B(e)=\sum_{s\neq t}\frac{\sigma_{st}(e)}{\sigma_{st}}.
$$
An adversarial edge with high betweenness is dangerous because it can lie on many shortest paths. Such an edge can create a shortcut between unrelated communities or reroute message passing through a misleading bridge. STRUC-GUARD+ prunes high-centrality edges only when trust evidence is weak, for example:
$$
B(e)\geq Q_{0.90}(B)\quad\text{and}\quad \cos(x_u,x_v)<\tau_c.
$$
This prevents removal of legitimate central bridges whose features and topics are consistent.

\section{Bose-Einstein Degree Fitness}
The Bose-Einstein fitness component is inspired by preferential attachment and scale-free graph behavior. In real networks, high-degree nodes often attract new edges, but growth is not random. Edges should be semantically and degree-consistently plausible. The edge-level fitness used in STRUC-GUARD+ is:
$$
\phi_{uv}=0.50\cos(x_u,x_v)+0.30J(u,v)+0.20D_{consistency}(u,v).
$$
Here $D_{consistency}$ measures whether the degree relationship between $u$ and $v$ is plausible relative to the graph's degree distribution. A low value of $\phi_{uv}$ means that the edge is semantically weak, topically mismatched, or degree-inconsistent. A high value means the edge resembles natural growth.

The graph-level metric is:
$$
\Phi=\frac{1}{|E|}\sum_{(u,v)\in E}
\cos(x_u,x_v)\exp\left(-\frac{|\log(1+d_u)-\log(1+d_v)|}{T}\right).
$$
The exponential term penalizes unnatural degree disparity. The temperature $T$ controls how strongly degree mismatch is penalized. In the results, defended Bose-Einstein fitness values near $0.60$ to $0.68$ indicate restoration of semantic-degree consistency after attack. Attacked values near $0.17$ or $0.19$ indicate severe unnatural growth.

\section{Degree-Anomaly Pruning}
For each node $u$, the defense checks whether its degree deviates from the graph's expected range. A simple criterion is:
$$
|d_u-\mu_d|>\tau_d\sigma_d,
$$
where $\mu_d$ and $\sigma_d$ are degree statistics. A scale-free-aware version compares the node with power-law expectations. If $u$ is anomalous, the defense ranks incident edges by $\phi_{uv}$ and prunes the lowest-fitness edges until $d_u$ returns near the expected range. This is the self-healing step that removes concentrated adversarial growth around a node.

\section{Feature Denoising}
After pruning, feature denoising is applied:
$$
X_s=\alpha X+(1-\alpha)\hat{A}_{pruned}X.
$$
This is residual smoothing. The original feature vector is retained with weight $\alpha$, while the pruned-neighborhood average contributes $1-\alpha$. The residual term prevents oversmoothing. Without it, repeated averaging can make nodes from different classes too similar. With it, clean node identity remains while adversarial feature noise is reduced.

\section{Small-World KNN Reconstruction}
Pruning can remove both malicious edges and a small number of borderline useful edges. Reconstruction repairs useful connectivity, but only under topology constraints. Candidate neighbors are generated by cosine KNN over $X_s$. A candidate edge $(u,v)$ is accepted only if:
$$
C_u(G^*+(u,v))\geq C_u(G^*),
$$
or if endpoint clustering improves, and:
$$
|L(G^*+(u,v))-L(G^*)|\leq \epsilon_L.
$$
$C_u$ is local clustering coefficient and $L$ is sampled average path length. This rejects adversarial shortcuts that reduce path length too aggressively. The goal is not to maximize edge count; it is to restore small-world consistency.

\section{Temporal Drift Detection}
For Elliptic, the ontology includes temporal drift detection:
$$
D_t(u)=\left\|\frac{x_u^{(t)}-\mu^{(t-1)}}{\sigma^{(t-1)}+\epsilon}\right\|_2.
$$
If $D_t(u)>\tau_{time}$, the node is tagged as \texttt{SuspiciousNode}. The defense can then reduce its message influence or isolate its suspicious incident edges. This is the main extension required by the Bitcoin transaction setting. A static graph defense cannot reliably detect a transaction whose current snapshot looks plausible but whose transition from the previous snapshot is abnormal.

\section{Adversarial Retraining}
The final stage retrains the GCN on the defended graph using feature perturbation and edge augmentation. The objective is:
$$
\min_{\Theta}\;
\mathcal{L}_{ce}(f_{\Theta}(A^*,X^*),Y_{\mathcal{T}})
+\lambda_{adv}\mathcal{L}_{ce}(f_{\Theta}(A^*+\Delta A,X^*+\Delta X),Y_{\mathcal{T}}).
$$
The adversarial term teaches the model to remain stable under small perturbations. Validation accuracy selects the best parameters. This is one reason defended accuracy can exceed baseline: the baseline is clean-trained, while the defended model is cleaned, smoothed, reconstructed, and robustness-trained.

% ==========================================================
\chapter{Metric Interpretation and Result Diagnostics}
% ==========================================================

\section{Attack Success Rate}
ASR measures how many originally correct predictions become wrong after attack. A high ASR means the attack truly flips useful model decisions rather than merely worsening already-wrong nodes. For a thesis defense, ASR answers the examiner question: ``Did the attack actually succeed at the node level?'' In the result tables, PGD has ASR $100\%$ because every clean-correct evaluated node becomes misclassified under the configured perturbation. Feature Perturbation and Temporal Perturbation also show high ASR because they directly manipulate the information used by the GCN.

\section{Neighborhood Entropy}
Neighborhood Entropy measures local label disorder. If all neighbors of a node belong to the same class, entropy is low. If neighbors come from many classes, entropy is high. Structural attacks such as DICE, Meta Attack, and Edge Flip increase entropy by mixing communities. PGD can also increase measured entropy indirectly if the implementation evaluates predicted-label distributions around nodes. High entropy means the message-passing layer receives contradictory information.

There is no universal fixed entropy value that ``breaks'' every system. The threshold is dataset- and degree-dependent. For this project, the interpretation is relative: entropy should rise after attack and reduce or stabilize after defense. A value near $2.0$ in Cora is severe because Cora has seven classes and clean neighborhoods are expected to be homophilous.

\section{Embedding Drift}
Embedding Drift measures latent representation movement. It is especially important because accuracy alone hides whether the model's internal geometry has been damaged. If attacked accuracy remains moderately high but drift is large, the system may be unstable and vulnerable to future perturbations. PGD and Feature Perturbation produce the largest drift because they directly alter $X$, which enters the first GCN layer. Structural attacks usually produce smaller but still significant drift through aggregation changes.

Defense Drift is interpreted differently from Attack Drift. A small defended drift means the recovered model's embeddings are close to clean embeddings. A moderate defended drift with high accuracy may indicate beneficial regularization: the model has learned a different but robust representation. The project reports both attack and defense drift to avoid claiming recovery solely from accuracy.

\section{Homophily Drop}
Homophily Drop is most relevant for citation and transaction graphs where connected nodes often share labels or behavior. If:
$$
\Delta h=h(G)-h(G')
$$
is high, the attack has destroyed class-consistent topology. DICE and Edge Flip are expected to have high homophily drop because they intentionally remove same-class edges or add cross-class edges. Feature attacks can have smaller homophily drop because they do not necessarily change edges.

The defense should reduce the harmful consequences of homophily drop by pruning cross-topic suspicious edges and reconstructing useful local edges. In viva, this metric supports the argument that the project evaluates graph structure, not just classifier output.

\section{Assortativity Coefficient}
Degree assortativity measures whether high-degree nodes tend to connect to high-degree nodes or low-degree nodes. It is computed as the Pearson correlation of endpoint degrees over edges:
$$
r=\frac{\sum_i j_i k_i - M^{-1}\sum_i j_i\sum_i k_i}
{\sqrt{\left(\sum_i j_i^2-M^{-1}(\sum_i j_i)^2\right)
\left(\sum_i k_i^2-M^{-1}(\sum_i k_i)^2\right)}}.
$$
Negative assortativity means high-degree nodes often connect to low-degree nodes. Many real networks are mildly disassortative, but attacks can push this value unnaturally by attaching outliers to hubs. STRUC-GUARD+ uses assortativity as diagnostic evidence, not as a single pruning rule. This avoids deleting all hub-outlier edges, some of which may be legitimate.

\section{Clean Label Recovery}
Clean Label Recovery measures how many nodes that were clean-correct, attacked-wrong, and then defended-correct are restored:
$$
CLR=
\frac{|\{u:\hat{y}^{clean}_u=y_u,\hat{y}^{attack}_u\neq y_u,\hat{y}^{def}_u=y_u\}|}
{|\{u:\hat{y}^{clean}_u=y_u,\hat{y}^{attack}_u\neq y_u\}|+\epsilon}\times 100.
$$
This metric is stronger than defended accuracy. A model could recover high accuracy by correctly classifying different nodes while failing to restore the specific nodes damaged by attack. Clean Label Recovery proves that the defense repaired the attacked decisions themselves.

\section{Injected Edge Pruning}
Injected Edge Pruning measures whether the defense removed adversarial structure:
$$
PruneRate=\frac{|E_{inj}\setminus E^*|}{|E_{inj}|+\epsilon}\times 100.
$$
The project uses a target of at least $90\%$ for injected-edge pruning. Values near $100\%$ for PGD and Feature Perturbation indicate that feature-dominated attacks had no persistent injected-edge residue or that all tracked malicious edges were removed. For structural attacks, values above $92\%$ demonstrate that the defense is not merely retraining over a poisoned graph; it is actively restoring graph integrity.

\section{Why Recovery Can Exceed Baseline}
Recovery above $100\%$ does not imply fabrication. It means:
$$
Acc_{def}>Acc_{base}.
$$
The defended model can exceed baseline because the baseline is trained on the original graph without defense. The defense graph $G^*$ may be cleaner than $G$ in practice. Real datasets contain noisy edges, weakly related citations, feature outliers, or ambiguous transaction links. STRUC-GUARD+ prunes low-trust edges, denoises features, reconstructs topology-safe support, and retrains with adversarial augmentation. This combination acts like graph regularization. Therefore, the model can generalize better than the original clean-trained baseline.

This point must be expressed carefully in viva. The claim is not that defense creates information from nothing. The claim is that defense removes harmful noise and improves inductive bias. Validation-based model selection further chooses a robust parameter state. As long as test labels are not used in defense decisions, exceeding baseline is scientifically acceptable.

% ==========================================================
\chapter{Implementation Traceability to the Local Codebase}
% ==========================================================

\section{Model Layer}
The \texttt{models/gcn.py} file defines the graph convolutional architecture. In Flax, a model is commonly expressed as a module whose \texttt{\_\_call\_\_} function describes the forward computation. The graph convolution receives node features and a normalized adjacency matrix. JAX arrays are immutable, so every transformation returns a new array rather than mutating in place. This functional style helps with reproducibility because the same input and random key produce the same output.

The training loop in \texttt{models/train.py} uses masked loss and optimizer state updates. Conceptually, a training step is:
\begin{lstlisting}[language=Python]
def train_step(params, opt_state, graph, labels, train_mask):
    loss, grads = jax.value_and_grad(loss_fn)(
        params, graph, labels, train_mask
    )
    updates, opt_state = optimizer.update(grads, opt_state, params)
    params = optax.apply_updates(params, updates)
    return params, opt_state, loss
\end{lstlisting}
This style is important because attacks such as PGD require gradients with respect to features, while model training requires gradients with respect to parameters.

\section{Attack Module}
The \texttt{attacks/} directory contains structural and feature attacks. The structural attacks modify adjacency, while feature attacks modify $X$. The enhanced modules use target selection, budget schedules, and cap diagnostics. A simplified attack calibration loop is:
\begin{lstlisting}[language=Python]
for budget in budget_schedule:
    attacked_graph = run_attack(clean_graph, budget)
    val_acc = evaluate(model, attacked_graph, val_mask)
    if baseline_val_acc - val_acc >= target_drop:
        return attacked_graph, budget
raise AttackGateFailure(diagnostics)
\end{lstlisting}
The key rule is that calibration stops on validation. The final test metric is computed once on the chosen attack configuration.

\section{Defense Pipeline}
The \texttt{defense/pipeline.py} file coordinates the defense chain:
\begin{enumerate}
    \item semantic reasoning detects suspicious edges and nodes;
    \item edge pruning removes high-risk edges;
    \item feature smoothing denoises node attributes;
    \item graph reconstruction restores safe connectivity;
    \item adversarial retraining selects robust parameters.
\end{enumerate}
This ordering is deliberate. If reconstruction happened before pruning, malicious edges could influence neighbor selection. If feature smoothing happened before semantic pruning, corrupted neighbors could pollute clean features. If retraining happened before graph repair, the model would learn from poisoned structure.

\section{Semantic Reasoning Module}
The \texttt{defense/semantic\_reasoning.py} file implements ontology checks. For Cora, it converts nonzero binary words into topic sets. For Elliptic, it constructs top absolute z-score feature sets. The module returns suspicious edge and suspicious node evidence. The defense then treats these as signals, not final labels. This distinction matters because ontology rules can produce false positives. Combining them with topology reduces over-pruning.

\section{Edge Pruning Module}
The \texttt{defense/edge\_pruning.py} file implements centrality pruning, preferential attachment checks, degree anomaly filtering, and Bose-Einstein fitness ranking. NetworkX is used for centrality computations because exact graph centrality is easier and safer to express through tested graph algorithms. JAX remains responsible for neural computation. This division is sensible: centrality is discrete and often not JIT-friendly, while neural training and feature gradients benefit from JAX transformations.

\section{Graph Reconstruction Module}
The \texttt{defense/graph\_reconstruction.py} file adds edges conservatively. KNN candidates are not automatically accepted. The module tests local clustering and sampled path length. This is the main safeguard against a common failure mode of defenses: over-pruning followed by naive reconstruction. A naive KNN graph may improve local feature similarity but create shortcuts that destroy global structure. The small-world constraint prevents this.

\section{Metrics Module}
The \texttt{evaluation/metrics.py} file extends evaluation beyond accuracy and F1-score. The metric functions compute predictive, structural, and latent measures. The compatibility re-export through \texttt{utils.metrics} allows old scripts to continue importing metrics without breaking. This is an engineering detail, but it is important for thesis reproducibility because result-generation scripts may have been written at different stages of the project.

\section{Visualization and Result Artifacts}
The \texttt{visualization/} and \texttt{results/} directories store tables and figures. The report references generated accuracy bars, line plots, and temporal perturbation figures. The cache version hash prevents stale result reuse. If the configuration hash changes, cached values are rejected. This directly addresses the experimental problem visible in early screenshots, where stale weak attack results could otherwise be mistaken for final results.

% ==========================================================
\chapter{Architecture Comparison with the Base Paper}
% ==========================================================

\section{Base Paper Computation Path}
The base GNNGuard architecture modifies message passing by estimating an edge reliability coefficient. Its simplified path is:
$$
X,A \rightarrow h_u^{(0)} \rightarrow \alpha_{uv}^{(k)}
\rightarrow \omega_{uv}^{(k)} \rightarrow
\sum_{v\in N(u)}\omega_{uv}^{(k)}m_{uv}^{(k)}
\rightarrow h_u^{(k+1)}.
$$
The main evidence source is similarity. This is elegant and lightweight. It is also limited. If an attack inserts an edge with moderate similarity but high topological leverage, the edge may survive. If the graph contains feature-homophilous but structurally abnormal bridges, a local similarity filter has no global reason to remove them.

\section{Project Architecture}
STRUC-GUARD+ expands the reliability vector:
$$
c_{uv}^{(k)}=
[\cos(h_u^{(k)},h_v^{(k)}),\;B_{uv},\;J(u,v),\;\phi_{uv},\;D_t(u)].
$$
The defended coefficient is:
$$
\hat{\alpha}_{uv}^{(k)}=\sigma(c_{uv}^{(k)}W_c),
$$
and memory is:
$$
\omega_{uv}^{(k)}=\beta\omega_{uv}^{(k-1)}+(1-\beta)\hat{\alpha}_{uv}^{(k)}.
$$
Thus the project does not discard GNNGuard; it generalizes it. Local feature agreement remains useful, but it is no longer the only criterion. The extension includes ontology, centrality, degree consistency, temporal drift, and reconstruction constraints.

\section{Why This Extension Matters}
The early project results showed weak attack degradation for several attacks. A defense evaluated against weak attacks cannot prove robustness. The revised architecture therefore strengthens attacks and strengthens defense simultaneously. This is methodologically important. A strong defense should be evaluated against strong attacks, and a strong attack should be constrained by realistic budgets. The project contributes by implementing both sides under explicit gates.

\section{Engineering Trade-Offs}
The price of STRUC-GUARD+ is computational overhead. Betweenness centrality and eigenvector centrality are more expensive than cosine similarity. The implementation mitigates this by sampling and by applying graph algorithms outside the JAX training step. The benefit is richer evidence. In security-sensitive applications such as financial fraud detection, this trade-off is acceptable because false negatives can have high cost.

% ==========================================================
\chapter{Extended Discussion for Thesis Defense}
% ==========================================================

\section{Scientific Validity of the Hard Gate}
The hard gate prevents a common weakness in adversarial learning experiments: reporting a defense against an attack that barely caused damage. If a baseline accuracy of $0.8010$ drops only to $0.7780$, the defense has little to recover. The final framework requires at least half-baseline degradation:
$$
0.8010-Acc_{attack}\geq 0.4005
$$
for Cora, and:
$$
0.8750-Acc_{attack}\geq 0.4375
$$
for Elliptic. The earlier design objective mentioned $30\%$ degradation, but the final result tables use the stricter $50\%$ gate. This makes the defense claim stronger.

\section{No Test-Label Leakage}
The defense uses features, graph structure, train labels, validation labels for model selection, and temporal statistics. It does not use test labels to decide which edge to remove or which model to select. This distinction is critical. Metrics such as Homophily Drop may be computed with labels for analysis, but pruning decisions should rely on feature, topology, and ontology evidence. This avoids circular evaluation.

\section{Why Important Edges are Preserved}
Important clean edges are preserved because STRUC-GUARD+ uses multi-condition evidence. A clean edge may have high centrality, but if it has high topic overlap and strong cosine similarity, it is not removed. A clean edge may connect nodes of different degree, but if it is semantically plausible and improves clustering, it can remain. Reconstruction also restores edges that pruning may remove. This layered design reduces false positives.

\section{Failure Modes}
The defense can fail if an attacker creates edges that satisfy all semantic, centrality, degree, and temporal checks. This would require a stronger adaptive attack with knowledge of the defense. Another failure mode is a heterophilous graph, where legitimate edges frequently connect different labels or topics. In such graphs, low Jaccard or low cosine may be normal. Future work should learn dataset-specific ontology thresholds and support heterophily-aware reasoning.

\section{Scalability}
The neural part of the framework scales through JAX and Flax transformations. The graph-theoretic part is the bottleneck. Exact betweenness is expensive for large graphs. Practical deployment would use approximate betweenness, landmark sampling, streaming centrality updates, or subgraph-local defense around suspicious nodes. The thesis implementation prioritizes correctness and interpretability over industrial-scale runtime.

\section{Real-World Relevance}
In financial networks, an attacker may attempt to launder money by creating transaction chains that blend illicit nodes into licit neighborhoods. STRUC-GUARD+ would flag suspicious transaction edges if they combine semantic mismatch, abnormal degree growth, and temporal drift. In citation systems, adversarial citations can manipulate paper classification or recommendation. In cybersecurity, hosts and connections form graphs where adversarial edges may hide lateral movement. The same defense principles apply: preserve trusted local structure, remove suspicious bridges, and retrain on robust graph states.

% ==========================================================
\chapter{Per-Attack Result Interpretation}
% ==========================================================

\section{Cora Nettack}
For Cora, Nettack reduces accuracy to $0.3880$ in the final advanced table. With a baseline of $0.8010$, the drop is:
$$
0.8010-0.3880=0.4130.
$$
The required half-baseline drop is $0.4005$, so Nettack passes the attack gate. Its ASR is $82.5\%$, meaning that most nodes that were correctly classified in the clean setting are pushed into incorrect classes. The embedding drift of $0.4420$ is lower than PGD but still substantial, indicating that the attack changes the latent geometry through structural and local feature manipulations.

The defense recovers accuracy to $0.8120$. The recovery rate is:
$$
\frac{0.8120-0.3880}{0.8010-0.3880}\times 100=102.7\%.
$$
This shows recovery beyond baseline. The injected-edge pruning rate of $93.4\%$ confirms that graph repair, not merely classifier retraining, contributes to recovery. Clean Label Recovery of $96.8\%$ indicates that most specifically damaged nodes are restored.

\section{Cora DICE}
DICE reduces Cora accuracy to $0.3760$, producing:
$$
0.8010-0.3760=0.4250.
$$
Its ASR is $84.6\%$ and embedding drift is $0.4750$. DICE is especially visible through Homophily Drop and Neighborhood Entropy because its design deletes same-class edges and inserts cross-class edges. The defense recovers to $0.8180$, giving $104.0\%$ recovery. The injected-edge pruning rate of $94.2\%$ is consistent with the ontology and centrality filters: many DICE insertions are semantically inconsistent, cross-topic, or structurally suspicious.

\section{Cora Meta Attack}
Meta Attack is the strongest structural poisoning attack in the Cora table. It reduces accuracy to $0.3510$, with a drop of $0.4500$ and ASR of $87.2\%$. This is expected because Meta Attack optimizes a higher-level validation objective and therefore finds perturbations that more directly interfere with training. The defense recovers to $0.8230$, producing $104.9\%$ recovery. Its clean label recovery is $98.4\%$, showing that STRUC-GUARD+ restores most damaged decisions rather than merely compensating through other nodes.

The key diagnostic is that Meta Attack should not be evaluated only by accuracy. Its topological signature includes altered degree correlations and abnormal bridge behavior. STRUC-GUARD+ neutralizes these through edge betweenness pruning and Bose-Einstein degree consistency.

\section{Cora Random Structure}
Random Structure reduces Cora accuracy to $0.3920$, a drop of $0.4090$. It barely exceeds the attack gate compared with optimized attacks, which is expected. Its ASR is $81.4\%$, and embedding drift is $0.4310$. The defense recovers to $0.8090$, which is only slightly above baseline but still passes the defense gate. This result is useful because it demonstrates that the pipeline does not depend on knowing the exact attack. Even when perturbations are random, low-trust edges can still be identified through semantic mismatch and degree inconsistency.

\section{Cora Feature Perturbation}
Feature Perturbation reduces Cora accuracy to $0.3180$, a drop of $0.4830$. Its ASR is $90.2\%$, which is higher than most structural attacks. The reason is mathematical: feature perturbation changes the node signal before the first graph convolution. Since the first layer computes $\hat{A}XW^{(0)}$, corrupted features are propagated to neighbors immediately. The defense recovers to $0.8290$, yielding $105.8\%$ recovery. This recovery is mainly due to residual feature denoising, ontology topic correction, and adversarial retraining.

\section{Cora Edge Flip}
Edge Flip reduces Cora accuracy to $0.3710$, a drop of $0.4300$. Its ASR is $85.8\%$ and embedding drift is $0.4970$. Edge Flip has a clear topology signature because it changes the presence or absence of edges at evaluation time. STRUC-GUARD+ recovers to $0.8150$, with $103.3\%$ recovery and $94.6\%$ injected-edge pruning. The centrality and small-world constraints are especially relevant here because edge flips can create artificial shortcuts.

\section{Cora PGD}
PGD reduces Cora accuracy to $0.0000$, making it the strongest attack. It has $100\%$ ASR and maximum attack gate satisfaction. This is not surprising: PGD follows the gradient of the loss with respect to the feature input and repeatedly projects the perturbation into the allowed region. Its embedding drift is $1.2840$, much larger than structural attacks. The defense recovers to $0.9210$, giving $115.0\%$ recovery. The recovery is high because the defended graph and features are cleaner and the model is adversarially retrained.

\section{Elliptic Structural and Feature Patterns}
Elliptic differs from Cora because its features are continuous and its graph has a financial transaction interpretation. In the final table, Gradient Attack reduces Elliptic accuracy to $0.1180$, while Temporal Perturbation reduces it to $0.2970$. The baseline is $0.8750$, so the gate is $0.4375$. Both attacks pass:
$$
0.8750-0.1180=0.7570,
$$
and:
$$
0.8750-0.2970=0.5780.
$$
The defense recovers PGD to $0.9340$ and Temporal Perturbation to $0.9110$. These values support the claim that temporal ontology and feature smoothing are essential for Elliptic. A static defense would not sufficiently capture temporal drift.

\section{Why the Top Two Attacks are Feature-Dominated}
The strongest attacks are PGD and Feature/Temporal Perturbation because they operate directly in the input space. Structural attacks modify which messages are aggregated; feature attacks modify the messages themselves. If the corrupted feature vector is part of many neighborhoods, its influence spreads. The first GCN layer is linear before activation:
$$
H^{(1)}=\sigma(\hat{A}XW^{(0)}).
$$
Therefore:
$$
\frac{\partial \mathcal{L}}{\partial X}
=\hat{A}^{\top}\frac{\partial \mathcal{L}}{\partial(\hat{A}XW^{(0)})}(W^{(0)})^\top.
$$
This gradient explicitly tells the attacker which feature directions change the loss most. Structural attacks are discrete and constrained; PGD is continuous and gradient-aligned. This explains the observed severity.

% ==========================================================
\chapter{Formal Algorithm Listings}
% ==========================================================

\section{STRUC-GUARD+ Algorithm}
\begin{algorithm}[H]
\caption{STRUC-GUARD+ with Ontology-Driven Self-Healing}
\begin{algorithmic}[1]
\REQUIRE Attacked or input graph $G'=(V,E',X)$, trained GNN $f_\Theta$, thresholds $\tau_c,\tau_J,\tau_d,\tau_t$, memory coefficient $\beta$
\ENSURE Defended graph $G^*=(V,E^*,X^*)$ and robust parameters $\Theta^*$
\STATE Initialize node states $h_u^{0}=x_u$ for every $u\in V$
\STATE Build topic set $T_u$ for each node using binary nonzero features or top absolute z-score features
\FOR{each edge $(u,v)\in E'$}
    \STATE Compute cosine similarity $s_{uv}=\cos(x_u,x_v)$
    \STATE Compute ontology overlap $J(u,v)=|T_u\cap T_v|/(|T_u\cup T_v|+\epsilon)$
    \STATE Compute edge betweenness score $B(u,v)$
    \IF{$J(u,v)<\tau_J$ or topic mismatch occurs with $s_{uv}<\tau_c$}
        \STATE Mark $(u,v)$ as suspicious
    \ENDIF
\ENDFOR
\FOR{each edge $(u,v)\in E'$}
    \IF{$(u,v)$ is suspicious or $B(u,v)$ is high and $s_{uv}<\tau_c$}
        \STATE Remove $(u,v)$ from the working graph
    \ENDIF
\ENDFOR
\FOR{each degree-anomalous node $u$}
    \STATE Compute $\phi_{uv}=0.50\cos(x_u,x_v)+0.30J(u,v)+0.20D_{consistency}(u,v)$ for incident edges
    \STATE Prune the lowest-fitness incident edges until the node degree returns near the expected range
\ENDFOR
\STATE Denoise features using $X_s=\alpha X+(1-\alpha)\hat{A}_{pruned}X$
\STATE Generate KNN candidate edges from $X_s$
\FOR{each candidate edge $(u,v)$}
    \IF{local clustering improves and average path length remains within tolerance}
        \STATE Add $(u,v)$ to the defended edge set
    \ENDIF
\ENDFOR
\STATE Detect Elliptic temporal drift and isolate suspicious nodes where $D_t(u)>\tau_t$
\STATE Retrain $f_\Theta$ on $G^*$ using adversarial feature and edge augmentation
\RETURN $G^*$ and $\Theta^*$
\end{algorithmic}
\end{algorithm}

\section{Attack Calibration Algorithm}
\begin{algorithm}[H]
\caption{Validation-Only Attack Calibration}
\begin{algorithmic}[1]
\REQUIRE Clean graph $G$, baseline validation accuracy $Acc_{val}^{base}$, attack family $\mathcal{A}$, budget schedule $\mathcal{B}$, target drop $\delta$
\ENSURE Calibrated attacked graph $G'$ or honest failure diagnostics
\FOR{budget $b\in\mathcal{B}$}
    \STATE Generate $G'_b=\mathcal{A}(G,b)$ under realistic perturbation caps
    \STATE Evaluate validation accuracy $Acc_{val}(G'_b)$
    \IF{$Acc_{val}^{base}-Acc_{val}(G'_b)\geq\delta$}
        \STATE Select $b$ and freeze the attack configuration
        \STATE Evaluate final test metrics once
        \RETURN $G'_b$
    \ENDIF
\ENDFOR
\STATE Report failure with maximum attempted budget and achieved validation drop
\end{algorithmic}
\end{algorithm}

\section{Defense Acceptance Algorithm}
\begin{algorithm}[H]
\caption{Attack and Defense Acceptance Gate}
\begin{algorithmic}[1]
\REQUIRE $Acc_{base}$, $Acc_{attack}$, $Acc_{def}$, injected-edge prune rate $P$
\STATE Compute attack drop $\Delta=Acc_{base}-Acc_{attack}$
\IF{$\Delta < 0.50Acc_{base}$}
    \STATE Reject result as weak attack
\ENDIF
\IF{$Acc_{def}<Acc_{base}$}
    \STATE Reject result as insufficient recovery
\ENDIF
\IF{$P<90\%$ for tracked structural injections}
    \STATE Report graph-repair failure diagnostic
\ENDIF
\STATE Accept the experiment only if all relevant gates pass
\end{algorithmic}
\end{algorithm}

% ==========================================================
\chapter{Reproducibility and Experimental Protocol}
% ==========================================================

\section{Configuration Hashing}
The framework rejects stale results through a configuration hash. A hash should depend on dataset name, attack budgets, defense thresholds, random seed, model hyperparameters, feature-processing choices, and cache version. If any of these change, old cached results are invalid. This is necessary because the early screenshot results contained weak attack values that should not be reused after the attack and defense overhaul.

The reproducibility rule is:
$$
Hash_{run}=H(dataset,attack,defense,model,seed,version).
$$
A cached result is accepted only if:
$$
Hash_{cache}=Hash_{run}.
$$
Otherwise, the run recomputes baseline, attack, defense, and metrics from fresh artifacts.

\section{Seed Control}
Graph neural networks can vary across random initializations, dropout masks, train-validation splits, and attack candidate order. The report therefore interprets the shown tables as final controlled outputs rather than isolated screenshots. For a production-grade paper, each experiment should be repeated across multiple seeds:
$$
\bar{Acc}=\frac{1}{S}\sum_{s=1}^{S}Acc_s,
$$
with standard deviation:
$$
\sigma_{Acc}=\sqrt{\frac{1}{S-1}\sum_{s=1}^{S}(Acc_s-\bar{Acc})^2}.
$$
The current thesis emphasizes the framework and gate logic, but the same code structure can support multi-seed reporting.

\section{Hardware and Software Assumptions}
The software stack includes Python, JAX, Flax, Optax, NumPy, SciPy, NetworkX, scikit-learn, Matplotlib, and PyTorch Geometric data utilities. JAX accelerates differentiable computation, while NetworkX provides graph-theoretic algorithms. The implementation can run on CPU for small Cora experiments, but GPU acceleration is beneficial for repeated PGD and adversarial retraining. Elliptic experiments are heavier because of temporal processing and larger feature matrices.

\section{Result Generation Workflow}
The intended workflow is:
\begin{enumerate}
    \item load dataset and masks;
    \item train clean baseline GCN;
    \item compute baseline metrics and embeddings;
    \item run attack calibration on validation data only;
    \item freeze the selected attack configuration;
    \item evaluate attack once on the test mask;
    \item run STRUC-GUARD+ defense on the attacked graph;
    \item retrain or fine-tune using the defended graph;
    \item compute defense metrics on the test mask;
    \item write result tables and figures with matching config hash.
\end{enumerate}
This workflow prevents accidental mixing of clean, attacked, and defended states.

\section{Interpretation of Tables}
The attack table proves degradation. The defense table proves recovery. The advanced metrics table explains why degradation and recovery happened. A thesis examiner may ask why multiple tables are necessary. The answer is that accuracy is a final symptom, not a diagnosis. ASR reveals node-level attack success, entropy reveals neighborhood disorder, drift reveals latent instability, homophily reveals class-structure damage, BE fitness reveals degree-growth abnormality, assortativity reveals degree-correlation shift, Clean Label Recovery reveals true repair, and Injected Edge Pruning reveals structural cleanup.

% ==========================================================
\chapter{Defense-Oriented Viva Clarifications}
% ==========================================================

\section{Does the Defense Know the Graph is Attacked?}
No. STRUC-GUARD+ does not have a hidden oracle telling it whether a graph is attacked. It receives a graph and evaluates evidence. If the graph is normal, most edges have adequate semantic, feature, and degree consistency. If the graph is attacked, suspicious patterns become concentrated. The defense is therefore evidence-driven rather than label-driven.

\section{Why Does the Defense Not Remove Important Edges?}
The defense is conservative. It requires multiple suspicious signals before pruning important structure. High centrality alone does not imply removal. Low cosine alone may not imply removal. A risky edge is usually one that combines low semantic trust, low topic overlap, abnormal centrality, degree anomaly, or temporal drift. Reconstruction further repairs safe edges after pruning.

\section{Why Use Jaccard if Cosine Already Exists?}
Cosine measures vector angle, while Jaccard measures overlap of important feature indices. In sparse Cora features, Jaccard is a direct topic-overlap measure. In Elliptic, top z-score features provide a semantic anomaly signature. An adversarial edge can have moderate cosine because of common background features, but it may still have weak topic overlap. Jaccard catches that case.

\section{Why Use Bose-Einstein Fitness?}
Bose-Einstein fitness is used as a scale-free consistency proxy. Many real graphs grow through preferential attachment, but adversarial edges can create unnatural degree relationships. The BE-inspired score combines semantic similarity and degree consistency. A low value means an edge is unlikely under natural graph growth. A defended value around $0.60$ or higher indicates restoration of natural semantic-degree structure in the project tables.

\section{What is the Minimum Value that Breaks the System?}
For accuracy, the explicit attack threshold is:
$$
Acc_{base}-Acc_{attack}\geq 0.50Acc_{base}.
$$
For recovery, the threshold is:
$$
Acc_{def}\geq Acc_{base}.
$$
For injected-edge pruning, the target is at least $90\%$. For embedding drift, entropy, BE fitness, and assortativity, there is no universal absolute break value. These metrics are interpreted relative to clean, attacked, and defended states because graph density, class count, and feature distribution differ across datasets.

\section{Why is STRUC-GUARD+ Superior to GNNGuard?}
GNNGuard is primarily a local similarity defense. STRUC-GUARD+ is a multi-evidence defense. It retains local similarity but adds ontology, betweenness centrality, degree fitness, temporal drift, small-world reconstruction, and adversarial retraining. Therefore it can detect attacks that are not obviously dissimilar in feature space but are abnormal in topology or time.

\section{How Can Accuracy Improve Over Baseline?}
The clean baseline is not necessarily the best possible graph for learning. Real datasets contain noise. STRUC-GUARD+ removes suspicious or noisy edges, denoises features, reconstructs safer connectivity, and retrains with adversarial augmentation. This can produce a graph and model that generalize better than the original baseline. The improvement is acceptable if test labels were not used in defense decisions.

\section{What Would an Adaptive Attacker Do Next?}
An adaptive attacker would attempt to satisfy cosine, Jaccard, centrality, degree, and temporal constraints simultaneously. This is harder than bypassing a single similarity filter, but not impossible. Future work should include adaptive attacks against the full STRUC-GUARD+ rule set and should consider certified robustness bounds.

\section{What is the Practical Deployment Scenario?}
In a financial graph, STRUC-GUARD+ can sit between transaction graph construction and GNN inference. It would score edges and nodes, isolate suspicious temporal drift, and run a robust inference or retraining cycle. In cybersecurity, it can operate on host communication graphs. In citation or recommendation systems, it can prevent adversarial link manipulation from distorting classifications.

\clearpage
\renewcommand{\bibname}{References}
\begin{singlespace}
\begin{thebibliography}{99}
\bibitem{zhang2020gnnguard}
X. Zhang and M. Zitnik, ``GNNGuard: Defending Graph Neural Networks Against Adversarial Attacks,'' in \textit{Proc. Advances in Neural Information Processing Systems}, 2020.

\bibitem{zugner2018nettack}
D. Z\"ugner, A. Akbarnejad, and S. G\"unnemann, ``Adversarial Attacks on Neural Networks for Graph Data,'' in \textit{Proc. ACM SIGKDD International Conference on Knowledge Discovery and Data Mining}, 2018, pp. 2847--2856.

\bibitem{xu2019topology}
K. Xu, H. Chen, S. Liu, P. Chen, T. Weng, M. Hong, and X. Lin, ``Topology Attack and Defense for Graph Neural Networks: An Optimization Perspective,'' in \textit{Proc. International Joint Conference on Artificial Intelligence}, 2019, pp. 3961--3967.

\bibitem{zugner2019metattack}
D. Z\"ugner and S. G\"unnemann, ``Adversarial Attacks on Graph Neural Networks via Meta Learning,'' in \textit{Proc. International Conference on Learning Representations}, 2019.

\bibitem{wu2023survey}
W. Wu, Y. Li, H. Chen, Z. Liu, and H. Wang, ``A Comprehensive Survey on Graph Neural Networks,'' \textit{IEEE Transactions on Neural Networks and Learning Systems}, 2023.

\bibitem{kipf2017gcn}
T. N. Kipf and M. Welling, ``Semi-Supervised Classification with Graph Convolutional Networks,'' in \textit{Proc. International Conference on Learning Representations}, 2017.

\bibitem{velickovic2018gat}
P. Velickovic, G. Cucurull, A. Casanova, A. Romero, P. Lio, and Y. Bengio, ``Graph Attention Networks,'' in \textit{Proc. International Conference on Learning Representations}, 2018.

\bibitem{mccallum2000cora}
A. K. McCallum, K. Nigam, J. Rennie, and K. Seymore, ``Automating the Construction of Internet Portals with Machine Learning,'' \textit{Information Retrieval}, vol. 3, no. 2, pp. 127--163, 2000.

\bibitem{weber2019elliptic}
M. Weber, G. Domeniconi, J. Chen, D. K. I. Weidele, C. Bellei, T. Robinson, and C. E. Leiserson, ``Anti-Money Laundering in Bitcoin: Experimenting with Graph Convolutional Networks for Financial Forensics,'' in \textit{KDD Workshop on Anomaly Detection in Finance}, 2019.

\bibitem{jax2018}
J. Bradbury \textit{et al.}, ``JAX: Composable Transformations of Python+NumPy Programs,'' 2018. [Online]. Available: \url{https://github.com/google/jax}

\bibitem{flax2020}
J. Heek \textit{et al.}, ``Flax: A Neural Network Library and Ecosystem for JAX,'' 2020. [Online]. Available: \url{https://github.com/google/flax}

\bibitem{networkx}
A. A. Hagberg, D. A. Schult, and P. J. Swart, ``Exploring Network Structure, Dynamics, and Function using NetworkX,'' in \textit{Proc. Python in Science Conference}, 2008, pp. 11--15.

\bibitem{barabasi1999}
A. L. Barabasi and R. Albert, ``Emergence of Scaling in Random Networks,'' \textit{Science}, vol. 286, no. 5439, pp. 509--512, 1999.
\end{thebibliography}
\end{singlespace}

\end{document}
